% Sections S1 to S6 of the arXiv version. They are input by
% problem131_arxiv.tex, which appends them to the paper. The journal version
% does not contain them (they formed Online Resource 1 until October 24, 2026).
% Edit the text here only.

\section{The center and the crossing to second order}\label{sec:erw-second}

Theorem~\ref{thm:erw-bounds} gives $\log C_N$ up to an error
$O(\varepsilon\log N)$. A lower bound that matches part (b) of that theorem
brings the error down to $O(\varepsilon)$.

\begin{proposition}[the center to second order]\label{prop:second-order}
Let $N\ge64$ be even. For every $0<\varepsilon<\frac12$,
\[
 -11\varepsilon-\varepsilon^2(84\log N+1.05\cdot10^5)
 \le\log C_N-\log\frac{4\varepsilon}{N+2}-2\varepsilon\log N .
\]
If moreover $\varepsilon(1+\log N)\le\frac1{100}$, then the middle term is at most
$5200\,\varepsilon$, and so
$\bigl\lvert\log C_N-\log\frac{4\varepsilon}{N+2}-2\varepsilon\log N\bigr\rvert
\le5200\,\varepsilon$.
\end{proposition}

So $C_N$ is $\frac{4\varepsilon}{N+2}N^{2\varepsilon}$ times a factor
$1+O(\varepsilon)$.

\paragraph{Idea of the proof.}
The proof is in Section~\ref{app:erw-bridge}. The upper bound is
Theorem~\ref{thm:erw-bounds}(b). For the lower bound we compare $M$ after
the first flip with P\'olya's urn conditioned to end at $h$, which we call
the bridge. Then $V_{t+1}(1)/\pi_t$ is the mean over the bridge of a product
of factors, one for each step (Proposition~\ref{prop:A-bridge}). Each factor
is $1$ plus $\varepsilon$ times a ratio that need not be bounded uniformly in
$N$, and we control the factors through their first two moments under the
bridge. By Jensen's inequality the mean of the
product is at least the exponential of the mean of the sum of the
logarithms. The step at time $u$ adds about $2\varepsilon/u$ to this mean,
and the sum over $u\le N$ is about $2\varepsilon\log N$.

With this bound the proof of Theorem~\ref{thm:erw-crossing} gives one more
term of the crossing.

\begin{proposition}[the crossing to second order]\label{prop:refined}
As $N\to\infty$ over even $N$,
\[
 x_N=g_N-\frac{2g_N\log N+g_N^2/2}{N(1+1/g_N)}+O\Bigl(\frac{\log N}N\Bigr).
\]
More precisely, for every even $N$ with $\log N\ge30$ the error term is at most
$20822\log N/N$ in absolute value. So $x_N<g_N$ for every even $N$ with $\log N\ge3500$.
\end{proposition}

So the root $g_N$ overshoots the crossing by an amount of order $\log^2N/N$. The
amount is $6\log^2N/N$ in the limit, but only about 60 percent of that at $N=51200$
(Table~\ref{tab:erw-crossing}).

\begin{proof}
We repeat the proof of Theorem~\ref{thm:erw-crossing}(a) with one more term
on each side. The center gains the term $2\varepsilon\log N$ and the end
gains the term $-x^2/(2N)$, and both push $F$ up.

\emph{Step 1. Both sides.} Let $F$ be as in the proof of
Theorem~\ref{thm:erw-crossing}, and let $1\le x\le2\Lambda_N$ and
$\varepsilon=x/N$. For large $N$ Proposition~\ref{prop:second-order} applies,
and its error $5200\,\varepsilon$ is $O(x/N)$, as is the term
$\log(1+2/N)$ in $\log\frac{4\varepsilon}{N+2}$. For the end we use
$\log(1-y)=-y-\frac{y^2}2+O(y^3)$ for $0\le y\le\frac12$ with $y=x/N$. So
\[
 \log C_N=\log x+\log4-2L+\frac{2xL}N+O\Bigl(\frac xN\Bigr),\qquad
 \log E_N=-\log2-x-\frac{x^2}{2N}+O\Bigl(\frac xN\Bigr).
\]
In the second formula we multiplied the expansion by $N-1$, and the terms
left over are $O(x/N)$ because $x^2/N$ is bounded for $x\le4L$. Subtracting
gives
\[
 F(x)=G(x)-\Lambda_N+\frac{2xL+x^2/2}{N}+O\Bigl(\frac xN\Bigr)
 \qquad(1\le x\le2\Lambda_N).
\]

\emph{Step 2. Solve for $x_N$.} By Theorem~\ref{thm:erw-crossing}, $x_N$ lies
in this range for large $N$. So $F(x_N)=0$ gives $G(x_N)=\Lambda_N-\eta_N$
with $\eta_N=(2x_NL+x_N^2/2)/N+O(L/N)$. Since $G(g_N)=\Lambda_N$, the mean
value theorem gives
$x_N-g_N=-\eta_N/G'(\xi)$ for some $\xi$ between $x_N$ and $g_N$. Here
$G'(\xi)=1+1/\xi=(1+1/g_N)(1+O(1/N))$, since $\lvert\xi-g_N\rvert=O(L^2/N)$
and $g_N\ge L$ for large $N$. As $\eta_N=O(L^2/N)$, we get
$x_N-g_N=-\eta_N/(1+1/g_N)+O(L^2/N^2)$. Replacing $x_N$ by $g_N$ in the main
term of $\eta_N$ changes it by $O(L^3/N^2)$, again because
$\lvert x_N-g_N\rvert=O(L^2/N)$. All these errors are $O(L/N)$, which gives
the formula. Its main term is of order $L^2/N$ and is larger than the error
for large $N$, so $x_N<g_N$.

\emph{Step 3. The explicit bound.} Let $L\ge30$. By the proof of
Theorem~\ref{thm:erw-crossing}, $24L^4+20800L^3\le N$, and $d=x_N-g_N$ has
$\lvert d\rvert\le\rho_N\le1/L^2$. Let $1\le x\le4L$ and $\varepsilon=x/N$. Then
$\varepsilon(1+\log N)\le\frac1{100}$, so Proposition~\ref{prop:second-order} bounds
$\log C_N-\log\frac{4\varepsilon}{N+2}-2\varepsilon L$ by $5200\,\varepsilon$ in absolute value.
Here $\log\frac{4\varepsilon}{N+2}=\log x+\log4-2L-\log(1+\frac2N)$ with
$0\le\log(1+\frac2N)\le\frac{2x}N$. For $0\le y\le\frac12$ the series of the logarithm gives
$\lvert\log(1-y)+y+\frac{y^2}2\rvert\le\frac{y^3}{3(1-y)}\le2y^3$. So
$(N-1)\log(1-\varepsilon)=-x+\varepsilon-\frac{x^2}{2N}+\frac{\varepsilon^2}2+R$ with
$\lvert R\rvert\le2N\varepsilon^3=2x^3/N^2\le x/N$, because $2x^2\le32L^2\le N$. Collecting
terms, and using $\frac{\varepsilon^2}2\le\varepsilon$,
\begin{equation}\tag{$\ast\ast$}
 \Bigl\lvert F(x)-G(x)+\Lambda_N-\frac{2xL+x^2/2}N\Bigr\rvert
 \le5200\varepsilon+\frac{2x}N+2\varepsilon+\frac xN=\frac{5205\,x}N .
\end{equation}
Now $x_N\in[1,4L]$ and $F(x_N)=0$. Put $\eta=(2g_NL+g_N^2/2)/N$. Since
$\Lambda_N/2\le\Lambda_N-\log\Lambda_N\le g_N\le\Lambda_N$, we have $L/2\le g_N\le2L$, so
$\lvert d/g_N\rvert\le\frac12$. Also $G(x_N)-\Lambda_N=d+\log(1+d/g_N)$, and
$d/g_N-(d/g_N)^2\le\log(1+d/g_N)\le d/g_N$. The main term of ($\ast\ast$) at $x_N$ differs from
$\eta$ by $d(2L+\frac{x_N+g_N}2)/N$, which is at most $5L\rho_N/N\le L/N$ in absolute value.
So ($\ast\ast$) at $x_N$ gives
\[
 \Bigl\lvert d\Bigl(1+\frac1{g_N}\Bigr)+\eta\Bigr\rvert\le\frac{5205\,x_N}N+\frac LN
 +\Bigl(\frac d{g_N}\Bigr)^2\le\frac{20822\,L}N ,
\]
since $5205\,x_N\le20820L$ and $(d/g_N)^2\le4\rho_N^2/L^2\le L/N$. The last inequality holds
because $\rho_N\le718L^2/N$ and $4\cdot718^2L\le20800L^3\le N$. Dividing by $1+1/g_N\ge1$ gives
$\lvert x_N-g_N+\eta/(1+1/g_N)\rvert\le20822L/N$, which is the explicit bound.

Finally let $L\ge3500$. The tangent of the logarithm at $8192$ gives
$\log\Lambda_N\le\Lambda_N/8192+\log8192-1$, so $g_N\ge\Lambda_N-\log\Lambda_N\ge1.996L$. Then
$2g_NL+g_N^2/2\ge5.984L^2$ and $1+1/g_N\le1+1/(1.996L)$. As
$20822L+20822/1.996<5.984L^2$ for $L\ge3500$, we get $\eta/(1+1/g_N)>20822L/N$, and the
explicit bound gives $x_N<g_N$.
\end{proof}

\begin{remark}[numerical values and open points]\label{rem:refined}
The formula of Proposition~\ref{prop:refined} without its error term is
within $0.07$ of $x_N$ already at $N=400$ (Table~\ref{tab:second}).
Proposition~\ref{prop:second-order} does not identify the term of order
$\varepsilon$. At fixed $N$ the center is a polynomial
$c_1\varepsilon+c_2\varepsilon^2+\cdots$ in $\varepsilon$, by
Proposition~\ref{prop:first-mutation}(b). From its Taylor coefficients,
computed in double precision (\texttt{verify\_elephant\_crossing.py}), we get
$c_2/c_1-2\log N=-1.585878$, $-1.546572$ and $-1.541005$ at $N=1600$, $12800$
and $51200$. These values approach $2\gamma-2-\log2=-1.5387\ldots$, but we do
not prove this limit here. We do not know whether, for some fixed $k$, the
bound
$\log C_N\le\log(c_1\varepsilon)+\varepsilon c_2/c_1+O(\varepsilon^2(\log N)^k)$
holds uniformly in $N$ in the range of Proposition~\ref{prop:second-order}.
Table~\ref{tab:erw-crossing} has $x_N<g_N$ already at $N=50$, while
Proposition~\ref{prop:refined} proves it only for $\log N\ge3500$.
\end{remark}

\begin{table}[htbp]
\centering\small
\caption{The crossings $x_N$ of Table~\ref{tab:erw-crossing}, their
distance to $x_{\mathrm{2nd}}$, the formula of
Proposition~\ref{prop:refined} without its error term, and the difference
$\Upsilon_N$ of Theorem~\ref{thm:erw-markov}. These values are
numerical and are not used in any proof.}
\label{tab:second}
\begin{tabular}{@{}rlll@{}}
\toprule
$N$ & $x_N$ & $x_N-x_{\mathrm{2nd}}$ & $\Upsilon_N\log N$\\
\midrule
50 & 3.940487 & $+0.3452$ & $-0.083$\\
400 & 7.629217 & $+0.0621$ & $-0.199$\\
1600 & 10.240565 & $+0.0179$ & $+0.215$\\
3200 & 11.546834 & $+0.0097$ & $+0.386$\\
12800 & 14.158845 & $+0.0028$ & $+0.607$\\
51200 & 16.778468 & $+0.0008$ & $+0.707$\\
\bottomrule
\end{tabular}
\end{table}

\subsection{The bridge and the lower bound}\label{app:erw-bridge}

Fix $1\le t\le h$, and start the P\'olya$(c)$ chain at state $1$ at time $t+1$.
By Lemma~\ref{lem:A-polya}(b), given that it is at $h$ at time $N$ it moves
up at time $u$ from $b$ with probability $\bar r_u(b)=\frac{h-b}{N-u}$,
whatever $c$ is. So its $h-1$ up steps are placed uniformly at random among
the $N-t-1$ steps after time $t+1$. We call this law of the path
$(M_u)_{t+1\le u\le N}$ the bridge, and we write $\Prob^{\mathrm{br}}_t$ and
$\E^{\mathrm{br}}_t$ for it. For $1\le b\le u-1$ let
\[
 \varpi^+_u(b)=1+\varepsilon\frac{u-2b}{b},\qquad
 \varpi^-_u(b)=1-\varepsilon\frac{u-2b}{u-b},
\]
\[
 \theta_u(b)=\frac{N(u-2b)^2}{2(N-u)\,b(u-b)},\qquad
 Q_u(b)=(u-2b)^2\Bigl[\frac{\bar r_u(b)}{b^2}+\frac{1-\bar r_u(b)}{(u-b)^2}\Bigr],
\]
and let $\varpi_u=\varpi^+_u(M_u)$ if $M_{u+1}=M_u+1$ and
$\varpi_u=\varpi^-_u(M_u)$ otherwise. Put
\[
 \mathcal A_t=\sum_{u=t+1}^{N-1}\E^{\mathrm{br}}_t\,\theta_u(M_u),\qquad
 \mathcal B_t=\sum_{u=t+1}^{N-1}\E^{\mathrm{br}}_t\,Q_u(M_u).
\]
Both are nonnegative, and neither depends on $\varepsilon$.

\begin{proposition}[bridge formula]\label{prop:A-bridge}
Let $0\le\varepsilon\le1$ and $1\le t\le h$. Then
\[
 V_{t+1}(1)=\pi_t\,\E^{\mathrm{br}}_t\Bigl[\prod_{u=t+1}^{N-1}\varpi_u\Bigr].
\]
If $0<\varepsilon\le\frac12$, then
$\log V_{t+1}(1)\ge\log\pi_t+\varepsilon \mathcal A_t-\varepsilon^2\mathcal B_t$.
\end{proposition}

\begin{proof}
We divide $V$ by the hitting probability of P\'olya's urn. The quotient
satisfies a recursion whose weights are the step probabilities of the
bridge, and each step leaves one factor $\varpi_u$.
For $u\ge t+1$ and $b\in\mathcal S_u$ we have $\Phi^0_u(b)>0$, and we put
$R_u(b)=V_u(b)/\Phi^0_u(b)$. By Lemma~\ref{lem:A-polya}(a) and (b) with
$c=0$, $\frac bu\Phi^0_{u+1}(b+1)=\bar r_u(b)\Phi^0_u(b)$ and
$(1-\frac bu)\Phi^0_{u+1}(b)=(1-\bar r_u(b))\Phi^0_u(b)$. Dividing
$V_u=\mathcal L_uV_{u+1}$ by $\Phi^0_u(b)$ gives
\[
 R_u(b)=(1-\bar r_u(b))\,\frac{1-r_u(b)}{1-b/u}\,R_{u+1}(b)
 +\bar r_u(b)\,\frac{r_u(b)}{b/u}\,R_{u+1}(b+1),
\]
where a term with an infeasible state is absent, and its bridge probability
is $0$ (if $b=h$ then $\bar r_u(b)=0$, and if $u-b=h$ then $\bar r_u(b)=1$).
Also $\frac{r_u(b)}{b/u}=\varpi^+_u(b)$ and
$\frac{1-r_u(b)}{1-b/u}=\varpi^-_u(b)$. So
$R_u(b)=\E^{\mathrm{br}}_t[\varpi_uR_{u+1}(M_{u+1})\mid M_u=b]$, and
$R_N(h)=1$. Iterating from $u=t+1$ and using $\Phi^0_{t+1}(1)=\pi_t$ gives
the formula.

Now let $0<\varepsilon\le\frac12$. We take the logarithm of the formula and
bound it from below in three steps. First, every factor is at least
$\frac12$. Indeed if $b\le u/2$ then $\varpi^+_u(b)\ge1$ and
$\varpi^-_u(b)-1\in[-\varepsilon,0]$, and if $b\ge u/2$ the same holds with
$+$ and $-$ exchanged. Second, by Jensen's inequality for the logarithm,
\[
 \log\E^{\mathrm{br}}_t\Bigl[\prod_u\varpi_u\Bigr]
 \ge\sum_u\E^{\mathrm{br}}_t\log\varpi_u
 \ge\sum_u\E^{\mathrm{br}}_t\bigl[(\varpi_u-1)-(\varpi_u-1)^2\bigr].
\]
The second inequality is $\log(1+y)\ge y-y^2$ for $y\ge-\frac12$, which holds
since the difference vanishes at $0$ and its derivative $y(1+2y)/(1+y)$ has
the sign of $y$. Third, we compute the two means given $M_u=b$. Write
$\bar r=\bar r_u(b)$. The bridge moves up with probability $\bar r$, and
then $\varpi_u-1=\varepsilon\frac{u-2b}b$. Otherwise
$\varpi_u-1=-\varepsilon\frac{u-2b}{u-b}$. Since
$u\bar r-b=\frac{N(u-2b)}{2(N-u)}$, this gives
\[
 \E^{\mathrm{br}}_t[\varpi_u-1\mid M_u=b]=\varepsilon(u-2b)\frac{u\bar r-b}{b(u-b)}
 =\varepsilon\,\theta_u(b),\qquad
 \E^{\mathrm{br}}_t[(\varpi_u-1)^2\mid M_u=b]=\varepsilon^2Q_u(b).
\]
So the sum above is $\varepsilon \mathcal A_t-\varepsilon^2\mathcal B_t$, which is the bound.
\end{proof}

\begin{lemma}[the weights $w_t$]\label{lem:A-w}
For each fixed $t$, $w_t\to t2^{-t-1}$ as $N\to\infty$. For every even
$N\ge2$, $\sum_tw_t(t+1)^3\le124$.
\end{lemma}

\begin{proof}
Each fraction in $\pi_t=\frac t{N-t}\prod_{i=1}^{t-1}\frac{h-i}{N-i}$ tends
to $\frac12$, so $N\pi_t\to t2^{1-t}$, and $Nc_1\to4$. This is the first
claim. For the second claim we write $w_t$ and $(t+1)^3$ through binomial
coefficients in $t$. By the proof of Lemma~\ref{lem:urn} (with $v=t+1$ and
$d=h$), $\pi_t$ is a constant times $\binom t1\binom{N-1-t}{h-1}$, and the
Chu--Vandermonde identity gives
$\sum_t\binom tj\binom{N-1-t}{h-1}=\binom N{h+j}$. With $j=1$ this gives
the constant, namely $w_t=\binom t1\binom{N-1-t}{h-1}/\binom N{h+1}$. Now
$t(t+1)^3=8\binom t1+38\binom t2+54\binom t3+24\binom t4$, because both sides
are polynomials of degree four that agree at $t=0,1,2,3,4$. So
$\sum_tw_t(t+1)^3$ is the sum of $8$, $38$, $54$ and $24$ times
$\binom N{h+j}/\binom N{h+1}$ for $j=1,2,3,4$. Since
$\binom N{h+j}\le\binom N{h+1}$ for $j\ge1$, it is at most
$8+38+54+24=124$.
\end{proof}

Under $\Prob^{\mathrm{br}}_t$ let $K=M_u-1$ and $K'=m-K$ be the numbers of
up steps and of other steps among the $m=u-t-1$ steps between times $t+1$
and $u$. So $K$ is hypergeometric, with $m$ draws from $N-t-1$ steps of
which $h-1$ are up steps. Let $\Omega_u=\{u/8\le M_u\le7u/8\}$.

\begin{lemma}\label{lem:A-AB}
Let $N\ge16$, $1\le t\le h$ and $t+1\le u\le N-1$.
\begin{enumerate}[(a)]
\item The second moment is
\[
 \E^{\mathrm{br}}_t(u-2M_u)^2=\frac{(t-1)^2(N-u)^2}{(N-t-1)^2}
 +\frac{4(u-t-1)(h-1)(h-t)(N-u)}{(N-t-1)^2(N-t-2)}\le(t-1)^2+u .
\]
\item If $t\le\frac{N+2}4$, then
$\mathcal A_t\ge2(H_{N-1}-H_t-1-\frac1t)\ge2(\log N-\log t-3)$.
\item If $t\le\frac N8$ and $u\ge8(t+1)$, then
$\Prob^{\mathrm{br}}_t(\Omega_u^c)\le2e^{-0.0943u}$.
\item If $t\le\frac N8$, then $\mathcal B_t\le84(t+1+\log N)+683(t+1)^3+19100$.
\end{enumerate}
\end{lemma}

\begin{proof}
(a) Here $u-2M_u=u-2-2K$ and $u-2=m+t-1$. Since $N-2=2(h-1)$, the mean
$\E K=\frac{m(h-1)}{N-t-1}$ gives
\[
 \E(u-2M_u)=\frac{(m+t-1)(N-t-1)-m(N-2)}{N-t-1}
 =\frac{(t-1)(N-t-1-m)}{N-t-1}=\frac{(t-1)(N-u)}{N-t-1}.
\]
The variance of a hypergeometric variable with $m$ draws from $N-t-1$ steps,
of which $h-1$ are up steps and $h-t$ are not, is
$\Var K=\frac{m(h-1)(h-t)(N-t-1-m)}{(N-t-1)^2(N-t-2)}$, and $N-t-1-m=N-u$. So
the second term is $4\Var K$, and the identity is
$\E(u-2M_u)^2=(\E(u-2M_u))^2+4\Var K$. For the bound, the first term is at
most $(t-1)^2$ because $N-u\le N-t-1$. In the second term
$4(h-1)(h-t)=(N-2)(N-2t)=(N-t-1)^2-(t-1)^2\le(N-t-1)^2$, so it is at most
$\frac{m(N-u)}{N-t-2}$. This is $0$ for $u=t+1$, and for $u\ge t+2$ it is at
most $u$, because $m\le u$ and $N-u\le N-t-2$.

(b) We bound $\theta_u$ from below by a multiple of $(u-2b)^2$ and use only
the variance term of (a). Since $b(u-b)\le u^2/4$,
$\theta_u(b)\ge\frac{2N(u-2b)^2}{(N-u)u^2}$. By (a),
$\E^{\mathrm{br}}_t(u-2M_u)^2\ge4\Var K$, and in $4\Var K$ the factor $N-u$
cancels. With $4(h-1)(h-t)=(N-2)(N-2t)$ this gives
$\mathcal A_t\ge\nu_t\sum_{u=t+1}^{N-1}\frac{u-t-1}{u^2}$ with
$\nu_t=\frac{2N(N-2)(N-2t)}{(N-t-1)^2(N-t-2)}$. We show that $\nu_t\ge2$. Here
$\nu_1=\frac{2N}{N-3}>2$, and for real $t\in[1,\frac{N+2}4]$,
$\frac{d}{dt}\log\nu_t=-\frac2{N-2t}+\frac2{N-t-1}+\frac1{N-t-2}
\ge\frac3{N-t-1}-\frac2{N-2t}\ge0$,
since the last inequality is equivalent to $4t\le N+2$. Next, the sum is
$H_{N-1}-H_t-(t+1)\sum_{u=t+1}^{N-1}u^{-2}$, which is at least
$H_{N-1}-H_t-\frac{t+1}t$. The sum is also nonnegative. So multiplying by
$\nu_t\ge2$ gives the first inequality, whether this lower bound is positive
or not. For the second inequality use
$H_{N-1}-H_t\ge\log\frac N{t+1}\ge\log N-\log t-\frac1t$.

(c) Here $m\ge\frac78u$, so $\frac u8\le\frac m7$. Since $t\ge1$,
$\E K\ge\frac m2$, and since $t\le\frac N8$,
$\E K'=\frac{m(h-t)}{N-t-1}\ge\frac{3m}8$. Hoeffding's inequality also holds
for sampling without replacement \cite[Section~6]{hoeffding1963}, so
$\Prob(K-\E K\le-\sigma m)\le e^{-2\sigma^2m}$ for $\sigma>0$, and the same
holds for $K'$. Now $\{M_u<\frac u8\}\subset\{K-\E K<-\frac{5m}{14}\}$, and
since $u-M_u=t+K'$, also
$\{u-M_u<\frac u8\}\subset\{K'-\E K'<-\frac{13m}{56}\}$. So
$\Prob(\Omega_u^c)\le e^{-2(5/14)^2m}+e^{-2(13/56)^2m}\le2e^{-0.10778m}
\le2e^{-0.0943u}$.

(d) We split each term $\E^{\mathrm{br}}_tQ_u(M_u)$ of $\mathcal B_t$ into the part on
$\Omega_u$ and the part off $\Omega_u$, and the three terms of the bound
come from the parts on $\Omega_u$, the parts off $\Omega_u$ with $u<8(t+1)$,
and the parts off $\Omega_u$ with $u\ge8(t+1)$. In all of them we use that
$\bar r(u-b)^2+(1-\bar r)b^2\le u^2$, so that
$Q_u(b)\le\frac{(u-2b)^2u^2}{b^2(u-b)^2}$. On $\Omega_u$ we have
$b(u-b)\ge\frac{7u^2}{64}$, so
$Q_u\le\frac{4096}{49}\frac{(u-2b)^2}{u^2}\le83.6\frac{(u-2b)^2}{u^2}$. By (a)
these parts add up to at most
$83.6\sum_{u>t}(\frac{(t-1)^2}{u^2}+\frac1u)\le83.6(t+H_{N-1})
\le84(t+1+\log N)$.
Off $\Omega_u$ we only have $b(u-b)\ge u-1\ge\frac u2$, which gives
$Q_u\le4(u-2b)^2\le4u^2$. For $u<8(t+1)$ we bound the probability of
$\Omega_u^c$ by $1$, and these parts add up to at most
$\sum_{u<8(t+1)}4u^2\le\frac43(8(t+1))^3\le683(t+1)^3$. For $u\ge8(t+1)$,
(c) bounds the part by $8u^2e^{-0.0943u}$, and with $q=e^{-0.0943}$ we have
$8\sum_{u\ge1}u^2q^u=8q(1+q)/(1-q)^3<19100$.
\end{proof}

\begin{proof}[Proof of Proposition~\ref{prop:second-order}]
The upper bound is Theorem~\ref{thm:erw-bounds}(b). For the lower bound let
$0<\varepsilon<\frac12$ and $t_0=\lfloor N/8\rfloor\ge8$. Every step below uses only this
range of $\varepsilon$. By Proposition~\ref{prop:first-mutation}(a),
$C_N=\varepsilon c_1\sum_tw_t(1-\varepsilon)^{t-1}V_{t+1}(1)/\pi_t$. This is
$\varepsilon c_1$ times a mean over the probability vector $(w_t)$. By
Jensen's inequality the logarithm of the mean is at least the mean of the
logarithms. For $t\le t_0$ we bound $\log(V_{t+1}(1)/\pi_t)$ from below by
Proposition~\ref{prop:A-bridge}. For $t>t_0$ we use only that it is $\ge0$,
which is Lemma~\ref{lem:A-polya}(d). This gives
\[
 \log C_N\ge\log(\varepsilon c_1)+\sum_tw_t(t-1)\log(1-\varepsilon)
 +\sum_{t\le t_0}w_t\bigl(\varepsilon \mathcal A_t-\varepsilon^2\mathcal B_t\bigr).
\]
\emph{The first sum.} By $\log(1-\varepsilon)\ge-\varepsilon-\varepsilon^2$
and Proposition~\ref{prop:first-mutation}(c), it is at least
$-(\varepsilon+\varepsilon^2)\sum_tw_t(t-1)\ge-2\varepsilon-2\varepsilon^2$.

\emph{The terms with $\mathcal A_t$.} Since $t_0\le\frac{N+2}4$,
Lemma~\ref{lem:A-AB}(b) gives $\mathcal A_t\ge2(\log N-\log t-3)$ for $t\le t_0$. Let
$w_\le=\sum_{t\le t_0}w_t$. By Jensen's inequality
$\sum_tw_t\log t\le\log\sum_tw_tt\le\log3$, so
\[
 \sum_{t\le t_0}w_t\,\varepsilon \mathcal A_t\ge2\varepsilon(\log N-3)w_\le-2\varepsilon\log3
 \ge2\varepsilon\log N-2\varepsilon(1-w_\le)\log N-6\varepsilon-2\varepsilon\log3 .
\]
By Proposition~\ref{prop:first-mutation}(c),
$1-w_\le\le2\sum_{t>t_0}t2^{-t}=2(t_0+2)2^{-t_0}$, and $N<8(t_0+1)$. As
$(t_0+2)2^{-t_0}\log(8t_0+8)$ decreases for $t_0\ge8$, we get
$2\varepsilon(1-w_\le)\log N\le4\cdot10\cdot2^{-8}\log72\,
\varepsilon<0.67\,\varepsilon$.

\emph{The terms with $\mathcal B_t$.} Here $t\le t_0\le\frac N8$, so Lemma~\ref{lem:A-AB}(d) applies. With
$\sum_tw_t(t+1)\le4$ (Proposition~\ref{prop:first-mutation}(c)) and
Lemma~\ref{lem:A-w} it gives
$\sum_{t\le t_0}w_t\mathcal B_t\le84(3+1+\log N)+683\cdot124+19100
=84\log N+104128$.

Collecting terms, and using $c_1=\frac4{N+2}$,
\[
 \log C_N-\log\frac{4\varepsilon}{N+2}-2\varepsilon\log N
 \ge-(2+0.67+6+2\log3)\,\varepsilon-\varepsilon^2(84\log N+104130),
\]
which gives the lower bound since $2+0.67+6+2\log3<11$. For the last claim,
the hypothesis on $\varepsilon$ gives
$\varepsilon^2(84\log N+1.05\cdot10^5)\le\varepsilon\,
\frac{84\log N+1.05\cdot10^5}{100(1+\log N)}$. The fraction decreases in
$N$, and its value at $N=64$ is below $205$. Since $11+205<5200$, the
two-sided bound follows.
\end{proof}

\subsection{Comparison with the persistent walk}\label{or:markov}

We compare with the persistent walk of the introduction of the paper, in which each
step repeats the previous one with probability $p$ and switches with
probability $q=1-p$. Here $p$ and $q$ refer to that walk. Put $T=Nq$. Its end is $p^{N-1}/2$,
its crossing is $p_N$, and $z_N=N(1-p_N)$. So $x=N\varepsilon$ corresponds to $T$, and the
crossing $x_N$ corresponds to $z_N$. We compare the two crossings through the
crossing equation of~\cite[Proposition~3.8]{paperA}, which says that for every
$N\ge2$
\begin{equation}\label{eq:markov-crossing}
 z_N+\tfrac12\log z_N=L+c_0+\frac1{8z_N}+O(L^{-2}),\qquad
 c_0=\tfrac12\log\tfrac\pi8,
\end{equation}
where the error term is at most $\frac1{4L^2}+\frac{L^2+L}N$ in absolute
value, and that $L/2\le z_N\le L$. This is the only result of~\cite{paperA}
that we use. Doubling \eqref{eq:markov-crossing}
without its last two terms and adding $\log2$ gives $2z_N+\log(2z_N)=\Lambda_N+\log(2\pi)$, so
$2z_N-\log(2\pi)$ nearly solves $G(x)=\Lambda_N$. What is left is of order
$1/L$.

\begin{theorem}\label{thm:erw-markov}
Let $z_N$ be the crossing value of $T$ for the persistent walk with $N$ steps, as
above. As $N\to\infty$ over even $N$,
\[
 \Upsilon_N:=x_N-\bigl(2z_N-\log(2\pi)\bigr)=\frac{c_*}{\log N}
 +O\Bigl(\frac{\log\log N}{(\log N)^2}\Bigr),\qquad
 c_*=\tfrac12\log(2\pi)-\tfrac14=0.668939\ldots
\]
\end{theorem}

\begin{proof}
Write $\ell=\log(2\pi)$ and $\bar z=2z_N$. We compute $G(\bar z-\ell)$ in three steps
and then compare it with $G(x_N)$.

\emph{Step 1. The size of $z_N$.} By the bounds after \eqref{eq:markov-crossing},
the error term there is $O(L^{-2})$ and $L/2\le z_N\le L$. So
$\log z_N=\log L+O(1)$ and $\frac1{8z_N}=O(1)$, and
\eqref{eq:markov-crossing} gives $z_N=L-\frac12\log L+O(1)$. Hence
$\bar z=2L-\log L+O(1)$, and
\[
 \frac1{8z_N}=\frac1{8L}+\frac{L-z_N}{8Lz_N}=\frac1{8L}+O\Bigl(\frac{\log L}{L^2}\Bigr),
 \qquad\frac1{\bar z}=\frac1{2L}+O\Bigl(\frac{\log L}{L^2}\Bigr).
\]

\emph{Step 2. The equation for $\bar z$.} We double \eqref{eq:markov-crossing} and
add $\log2$. The left side becomes $2z_N+\log z_N+\log2=\bar z+\log\bar z$. On the
right, $2c_0+\log2=\log\frac\pi4$, and $\frac1{4z_N}=\frac1{4L}+O(\log L/L^2)$
by Step~1. So $\bar z+\log\bar z=2L+\log\frac\pi4+\frac1{4L}+O(\log L/L^2)$.

\emph{Step 3. The value $G(\bar z-\ell)$.} We have
$\log(\bar z-\ell)=\log\bar z+\log(1-\ell/\bar z)$, and
$\log(1-\ell/\bar z)=-\ell/(2L)+O(\log L/L^2)$ by Step~1. With Step~2 and
$\log\frac\pi4-\ell=-\log8$ this gives
\[
 G(\bar z-\ell)=\bar z+\log\bar z-\ell+\log\Bigl(1-\frac\ell{\bar z}\Bigr)
 =\Lambda_N+\frac1{4L}-\frac\ell{2L}+O\Bigl(\frac{\log L}{L^2}\Bigr)
 =\Lambda_N-\frac{c_*}L+O\Bigl(\frac{\log L}{L^2}\Bigr).
\]

\emph{Step 4. Compare with $x_N$.} By Theorem~\ref{thm:erw-crossing}(a),
$G(x_N)=\Lambda_N+O(L^2/N)$. So
$G(x_N)-G(\bar z-\ell)=c_*/L+O(\log L/L^2)$. Both $x_N$ and
$\bar z-\ell=2z_N-\log(2\pi)$ are $2L+O(\log L)$, so $G'(x)=1+1/x=1+O(1/L)$
between them. By the mean value theorem we may divide by $G'$, and we get
$x_N-(\bar z-\ell)=c_*/L+O(\log L/L^2)$.
\end{proof}

We computed $z_N$ from the exact law of the persistent walk~\cite{paperA}, with a tolerance
of about $10^{-15}$. The difference $\Upsilon_N$ changes sign between $N=400$ and $N=1600$, and
$\Upsilon_N\log N$ is $0.707$ at $N=51200$ against the limit $c_*=0.669$
(Table~\ref{tab:second}). So the asymptotic regime is reached slowly.

\section{The tail}\label{or:tail}

We keep the conventions of Section~\ref{sec:shape} and
Appendix~\ref{app:shape} of the paper, and we use the lemmas of that
appendix.

This section shows that the law falls like $x/(2m^2)$ beyond the humps.

\begin{theorem}[the tail]\label{thm:profile}
Let $N$ be even, and let $x=N\varepsilon\ge1$ depend on $N$ with
$x\log^2N/N\to0$. Let $\omega_N\to\infty$ and $\delta'_N\to0$. Then
\[
 f_N(m)=\frac{x}{m(m+1)}\bigl(1+o(1)\bigr)\qquad\text{and}\qquad
 P_N(N-m)=\frac{x}{2m^2}\bigl(1+o(1)\bigr),
\]
uniformly for $\omega_Nx\log^2(x+1)\le m\le\delta'_NN$. We do not compute
the absolute constants behind the $o(1)$ terms.
\end{theorem}

For $m$ well beyond $x\log^2x$ and well below $N$, the chance of landing in the
value at distance $m$ from an end is about half the expected number of flipped families with $m$
steps, namely $x/(2m^2)$. The half is there because the first step must point to that end.

\paragraph{Idea of the proof.}
The proof is below. For the lower bound we use the
first flipped family, which has size $m$ with probability about
$x/(m(m+1))$ (Lemma~\ref{lem:B-D}), while the other flips change the count
only by a relative amount of order $x\log m/m$. For the upper bound we show
$\Prob(M_N\ge m)\le\frac xm(1+O(x\log^2m/m))$, and the unimodality of $f_N$
turns this into a local bound. The mirror term $f_N(N-m)$ is at most
$C_N=O(x/N^2)$, which gives the constant $\frac12$.

\begin{lemma}[late start]\label{lem:B-late}
Let $j\ge2$, $t\ge j-1$ and $z\ge0$. Then $\Prob(M^{[j-1]}_t\ge
z)\le\exp\{-\frac{j-1}{2}(\frac zt-\varepsilon\log\frac{2t}{j-1})\}$.
\end{lemma}

\begin{proof}
This is a Chernoff bound with an exponent that changes with time. A chain
that starts at $0$ at the late time $j-1$ is like an urn that already holds
$j-1$ balls of one color, so the fraction $M^{[j-1]}_k/k$ moves slowly. For
this reason we take an exponent $\theta_k$ of size about $(j-1)/(2k)$.

Write $G_k=M^{[j-1]}_k$ and $\nu=(j-1)/2$. For $k\ge j-1$ let
$\theta_k=\log\frac k{k-\nu}$, so that $e^{\theta_k}-1=\frac\nu{k-\nu}$. We
claim that for $k\ge j-1$
\[
 \E e^{\theta_{k+1}G_{k+1}}\le e^{\varepsilon\nu/(k+1-\nu)}\,\E e^{\theta_kG_k}.
\]
To see this, condition on $G_k=c$. The chain moves up with probability
$r_k(c)\le\varepsilon+c/k$, and $1+y\le e^y$. So
\[
 \E\bigl[e^{\theta_{k+1}G_{k+1}}\bigm| G_k=c\bigr]
 =e^{\theta_{k+1}c}\bigl(1+r_k(c)(e^{\theta_{k+1}}-1)\bigr)
 \le\exp\Bigl\{c\Bigl(\theta_{k+1}+\frac{\nu}{k(k+1-\nu)}\Bigr)+\frac{\varepsilon\nu}{k+1-\nu}\Bigr\}.
\]
It remains to check that $\theta_{k+1}+\frac{\nu}{k(k+1-\nu)}\le\theta_k$.
Now $\theta_k-\theta_{k+1}=\log(1+w)$ with $w=\frac\nu{(k-\nu)(k+1)}$, and
$\log(1+w)\ge\frac w{1+w}=\frac\nu{k(k+1-\nu)}$. This proves the claim.

We start from $G_{j-1}=0$ and use the claim for $k=j-1,\dots,t-1$. Since
$\sum_{k=j}^t\frac1{k-\nu}\le\log\frac{t-\nu}\nu\le\log\frac{2t}{j-1}$, we
get $\E e^{\theta_tG_t}\le\exp\{\varepsilon\nu\log\frac{2t}{j-1}\}$. Finally
$\theta_t=-\log(1-\frac\nu t)\ge\frac\nu t$, so Markov's inequality gives
$\Prob(G_t\ge z)\le e^{-\nu z/t}\,\E e^{\theta_tG_t}$, which is the claim.
\end{proof}

\begin{proposition}[cumulative upper bound]\label{thm:B-cum-upper}
Assume~\eqref{eq:H}. For $6\le m\le N/2$ with
$m\ge16x\log m$ we have $\Prob(M_N\ge m)\le\frac
xm\bigl(1+800\,\frac{x\log^2m}m\bigr)$.
\end{proposition}

\paragraph{Idea of the proof.}
Given the first flip at time $j$, the count $M_N$ is at most $\lvert\mathcal T_V\rvert+Z_j$, the
size of the first flipped family plus the count from the rest of the tree.
The event $\lvert\mathcal T_V\rvert\ge m$ gives the main term $x/m$. The other terms come from a
shift of $\lvert\mathcal T_V\rvert$ by $Z_j$, from a large $Z_j$, and from a late first flip, and
each is smaller by a factor of order $x\log^2m/m$.

\begin{proof}
\emph{Step 1. Split by the first flip.} Let $J=\lfloor N/4\rfloor$ and
$Z_j=Z^{(j)}_{N-1}$. First let $2\le j\le J$. By Lemma~\ref{lem:firstmark},
given $V=j$ the count $M_N$ has the law of
$\lvert\mathcal T_V\rvert-Y_{\lvert\mathcal T_V\rvert}+Z^{(j)}_{N-\lvert\mathcal T_V\rvert}$, where,
given $V=j$, $\lvert\mathcal T_V\rvert$ has the law $\beta_j$ and is independent of $Y$ and
$Z^{(j)}$. This is at most $\lvert\mathcal T_V\rvert+Z_j$, because $Y\ge0$ and $Z^{(j)}$ is
nondecreasing in time. So $M_N\ge m$ needs $Z_j>m/8$, or $Z_j\le m/8$ and
$\lvert\mathcal T_V\rvert\ge m-Z_j$. Hence
\[
 \Prob(M_N\ge m\mid V=j)\le\Prob\Bigl(Z_j>\frac m8\Bigr)
 +\E\Bigl[\bar\beta_j(m-Z_j);\,Z_j\le\frac m8\Bigr].
\]
We compare $\bar\beta_j(m-z)$ with $\bar\beta_j(m)$ for an integer
$0\le z\le m/8$. By \eqref{eq:B-ratio}, each step down from $s$ to $s-1$
with $s\le m$ multiplies $\bar\beta_j$ by at most
$\varrho_j:=1+\frac{j-1}{N-m-j+2}$, and $\varrho_j\le1+\frac{4(j-1)}N$
because $N-m-j\ge N/4$. Since $\varrho^z-1\le z(\varrho-1)\varrho^z$ for
$\varrho\ge1$, we get
$\bar\beta_j(m-z)\le\bar\beta_j(m)\varrho_j^z\le\bar\beta_j(m)\bigl(1+z(\varrho_j-1)\varrho_j^{m/8}\bigr)$.
So the mean above is at most
$\bar\beta_j(m)\bigl(1+(\varrho_j-1)\varrho_j^{m/8}\,\E Z_j\bigr)$.

On $\{V>J\}$ there is no mark on $2,\dots,J$, so by Lemma~\ref{lem:tree}
$M_N$ is a copy of $M^{[J]}_N$ independent of $\{V>J\}$. Altogether
\[
 \Prob(M_N\ge m)\le\sum_{j=2}^J\Prob(V=j)\Bigl[\bar\beta_j(m)\bigl(1+(\varrho_j-1)\varrho_j^{m/8}\,\E Z_j\bigr)+\Prob\Bigl(Z_j>\frac m8\Bigr)\Bigr]
 +\Prob(V>J)\,\Prob\bigl(M^{[J]}_N\ge m\bigr).
\]
This gives four terms, which we bound in turn. The main term is
$\sum_j\Prob(V=j)\bar\beta_j(m)\le\Prob(\lvert\mathcal T_V\rvert\ge m)\le\frac xm$ by
Lemma~\ref{lem:B-D}.

\emph{Step 2. The shift term.} This is the sum over $2\le j\le J$ of
$\Prob(V=j)\,\bar\beta_j(m)\,(\varrho_j-1)\,\varrho_j^{m/8}\,\E Z_j$. We
bound the five factors by
\[
 \Prob(V=j)\le\varepsilon,\quad
 \bar\beta_j(m)\le e^{-\frac{(j-1)(m-1)}{N-1}},\quad
 \varrho_j-1\le\frac{4(j-1)}N,\quad
 \varrho_j^{m/8}\le e^{\frac{(j-1)m}{2N}},\quad
 \E Z_j\le x\log\frac N{j-1}.
\]
The second is \eqref{eq:B-subtail}, the fourth follows from the third and
$1+y\le e^y$, and the last is Lemma~\ref{lem:B-tail}. The two exponentials
combine. Since $m\ge6$ we have
$\frac{m-1}{N-1}-\frac m{2N}\ge\frac{m-2}{2N}\ge\frac m{3N}$. So with
$a=\frac m{3N}$ and $i=j-1$ the shift term is at most
\[
 \frac{4\varepsilon x}N\sum_{1\le i<N/4}i\,e^{-ai}\log\frac Ni .
\]
Here $i\le J-1<N/4$, so every term is nonnegative. In this sum only the first flips up to a time of order $1/a=3N/m$ matter,
because a later first flip rarely has a family of size $m$. To bound the
sum we write $\log\frac Ni=\log\frac m3+\log\frac1{ai}$. The first part is
nonnegative, and $\sum_{i\ge1}ie^{-ai}\le a^{-2}$. The second part is
positive only for $i<1/a$. There $i\log\frac1{ai}\le\frac1{ea}$, so these
at most $1/a$ terms add up to at most $\frac1{ea^2}$. Hence the shift term
is at most
\[
 \frac{4\varepsilon x}N\cdot\frac{\log m+1}{a^2}=\frac{36x^2(\log m+1)}{m^2} .
\]

\emph{Step 3. The big-$Z$ term.} This is
$\sum_{j=2}^J\Prob(V=j)\Prob(Z_j>m/8)$, and again
$\Prob(V=j)\le\varepsilon$. Recall that $Z_j$ counts the disagreeing steps
of a walk whose first $j-1$ steps agree. For small $j$ we only have the
general tail bound of Lemma~\ref{lem:B-tail}, which is of order
$x\log m/m$. For large $j$ the late start makes $Z_j\ge m/8$ exponentially
unlikely (Lemma~\ref{lem:B-late}). We split at
$j_{\mathcal K}=\lfloor\mathcal KN/m\rfloor+1$ with
$\mathcal K=32\log(m/x)$. Note that $m\ge16x\log6>28x$, so $\mathcal K\ge2$.

Let $2\le j\le j_{\mathcal K}$. Lemma~\ref{lem:B-tail} at the integer
$\lfloor m/8\rfloor+1$ gives $\Prob(Z_j>m/8)\le8xH_m/m$. There are at most
$\mathcal KN/m$ such $j$, and $\mathcal K\le32\log m$. So they contribute at
most
\[
 \varepsilon\cdot\frac{\mathcal KN}m\cdot\frac{8xH_m}m
 =\frac{8\mathcal Kx^2H_m}{m^2} .
\]

Let $j>j_{\mathcal K}$, so that $j-1>\mathcal KN/m$. We use
Lemma~\ref{lem:B-late} with $t=N-1$ and the level $m/8$. Its exponent
contains the term $\varepsilon\log\frac{2(N-1)}{j-1}$. Here
$\frac{2(N-1)}{j-1}\le\frac{2m}{\mathcal K}\le m$, so this term is at most
$\frac{x\log m}N\le\frac m{16N}$, by the hypothesis $m\ge16x\log m$. The
other term in the exponent is $\frac m{8(N-1)}\ge\frac m{8N}$. So the lemma
gives $\Prob(Z_j\ge m/8)\le e^{-(j-1)m/(32N)}$, which is a geometric
sequence in $j$. Summing over $j-1\ge j_{\mathcal K}>\mathcal KN/m$, these
$j$ contribute at most
\[
 \varepsilon\,\frac{e^{-\mathcal K/32}}{1-e^{-m/(32N)}}
 \le\varepsilon\cdot\frac xm\Bigl(\frac{32N}m+1\Bigr)\le\frac xm\cdot\frac{33x}m .
\]
Here $e^{-\mathcal K/32}=x/m$ by the choice of $\mathcal K$, and we used
$\frac1{1-e^{-y}}\le\frac1y+1$ for $y>0$ and $\varepsilon\le x/m$.

\emph{Step 4. No early flip.} If there is no flip up to time $J$, the walk
agrees with $X_1$ for about $N/4$ steps, and Lemma~\ref{lem:B-late} makes
$M_N\ge m$ exponentially unlikely in $m$. Namely
$\Prob(V>J)=(1-\varepsilon)^{J-1}\le e^{\varepsilon-3x/16}$, since
$J\ge3N/16$, and Lemma~\ref{lem:B-late} with $j-1=J\in[\frac{3N}{16},\frac N4]$
and $t=N$ gives $\Prob(M^{[J]}_N\ge m)\le\exp\{-\frac{3m}{32}+\frac
x8\log\frac{32}3\}$. As $x\le m/28$ and $\varepsilon\le\frac14$, the product
of the two is at most $e^{1/4-m/13}$.

\emph{Step 5. Collect.} Put $\ell=\log m\ge\log6>1.79$, so that $\ell^2>3.2$. We compare
each of the four error terms with $x^2\ell^2/m^2$. First $36(\ell+1)\le40\ell^2$. Next
$\mathcal K\le32\ell$ because $x\ge1$, and $H_m\le1+\ell\le1.6\,\ell$, so
$8\mathcal KH_m\le410\ell^2$. Also $33\le12\ell^2$. Finally $y^2e^{-y}\le4$ for $y\ge0$, which
with $y=m/13$ gives $e^{1/4-m/13}\le676e^{1/4}/m^2\le304\ell^2x^2/m^2$, since $x\ge1$ and
$676e^{1/4}<870<304\cdot3.2$. As $40+410+12+304\le800$, the four error terms add up to at most
$800x^2\ell^2/m^2$, which is the claim.
\end{proof}

\begin{proposition}[local upper bound]\label{thm:B-local-upper}
Let $N\ge64$ and $x\ge1$, and put $\bar c=\frac1{40}$ and $A_1=34000$. If $\varepsilon\log
N\le\bar c$ and $A_1x\log^2(x+1)\le m\le N/8$, then
\[
 f_N(m)\le\frac{x}{m(m+1)}\bigl(1+20000\,\Delta(m)^{1/2}\bigr),\qquad
 \Delta(m)=\frac{x\log^2m}m+\varepsilon\log N+\frac mN .
\]
\end{proposition}

\paragraph{Idea of the proof.}
To the right of its mode $f_N$ is nonincreasing, so $(k+1)f_N(m)$ is at most
the probability of $\{m-k\le M_N\le m\}$, which is a difference of two tails.
We bound the first tail from above by
Proposition~\ref{thm:B-cum-upper} and the second from below by summing
Proposition~\ref{thm:B-local-lower}. Then we choose $k$ to balance the
errors.

\begin{proof}
\emph{Step 0. Three facts.} Put $L_1=\log(x+1)\ge\log2$, so that $xL_1^2\ge0.48$,
$m\ge16200$ and $\log m\ge9$. Also $\log N\ge4$, so $\varepsilon\le\frac1{160}$. For
$0<A\le59000$ we have $\log(AxL_1^2)\le11+3L_1$, since $\log A\le11$, $\log x\le L_1$ and
$\log L_1\le L_1$. First, $\log(x+2)\le2L_1\le2.9L_1^2$, so
$804x\log(x+2)\le8A_1xL_1^2\le8m\le N$, and Lemma~\ref{lem:B-mode} gives
$m_h\le15400xL_1^2\le m/2$. Second, let $m'\ge m/2\ge s_0:=17000xL_1^2$. As $\log s/s$
decreases for $s\ge e$, we get
$16x\log m'/m'\le16x(11+3L_1)/s_0\le1$, because $16(11+3L_1)\le17000L_1^2$. So
$m'\ge16x\log m'$, and also $m'\ge6$. Third, $\log^2s/s$ decreases for $s\ge e^2$. At
$s_1=A_1xL_1^2$ we have $\log^2s_1\le(11+3L_1)^2\le242+18L_1^2$, and
$x(242+18L_1^2)\le s_1$. So $x\log^2m/m\le1$, and $\Delta(m)\le1+\frac1{40}+\frac18\le1.15$.

\emph{Step 1. A difference of tails.} Let $0\le k\le m/2$. Then
$m-k\ge m/2\ge m_h$, and
Lemma~\ref{lem:B-sandwich} with $a=m-k$ and $b=m$ gives
$(k+1)f_N(m)\le\Prob(M_N\ge m-k)-\Prob(M_N\ge m+1)$.

\emph{Step 2. Bound each tail.} Write $\Delta=\Delta(m)$. Here \eqref{eq:H} holds.
Proposition~\ref{thm:B-cum-upper} at $m-k$, which applies by Step~0, bounds the first term by
\[
 \frac x{m-k}+\frac{800x^2\log^2(m-k)}{(m-k)^2}\le\frac x{m-k}+\frac{3200x^2\log^2m}{m^2}
 \le\frac x{m-k}+\frac{3200x\Delta}m .
\]
For the second we sum Proposition~\ref{thm:B-local-lower} over $m<\ell\le\lfloor N/4\rfloor$.
The main terms give
$x\bigl(\frac1{m+1}-\frac1{\lfloor N/4\rfloor+1}\bigr)\ge\frac x{m+1}\bigl(1-\frac{8m}N\bigr)$.
As $\frac{\log(\ell+1)}\ell$ decreases for $\ell\ge3$ and $\sum_{\ell>m}\frac1{\ell(\ell+1)}\le\frac1{m+1}$,
the terms with $\frac{x\log(\ell+1)}\ell$ cost at most
$\frac x{m+1}\cdot20x\frac{\log(m+2)}{m+1}\le\frac x{m+1}\cdot\frac{40x\log^2m}m$. The terms
with $\varepsilon\log N$ cost at most $\frac x{m+1}\cdot20\varepsilon\log N$. The last terms are
$\frac{20x}{\ell(\ell+1)}\ell e^{-\ell/9}\le\frac{20x}{m+1}e^{-\ell/9}$, and they cost at most
$\frac{20x}{m+1}\cdot\frac{e^{-m/9}}{e^{1/9}-1}\le\frac x{m+1}\cdot200e^{-m/9}$. Since
$e^{-y}\le2/y^2$ for $y>0$ and $m\ge16200$, we have
$200e^{-m/9}\le32400/m^2\le2/m\le2x\log^2m/m$. So
$\Prob(M_N\ge m+1)\ge\frac x{m+1}(1-42\Delta)$. As $\frac x{m-k}-\frac
x{m+1}=\frac{x(k+1)}{(m-k)(m+1)}$ and $\frac{42x\Delta}{m+1}\le\frac{42x\Delta}m$, this gives
\[
 f_N(m)\le\frac{x}{(m-k)(m+1)}+\frac{3242\,x\Delta}{(k+1)m} .
\]

\emph{Step 3. Choose $k$.} Put $\delta=\Delta^{1/2}$ and $X=\frac x{m(m+1)}$. Note that
$\frac1m\le\frac{x\log^2m}m\le\Delta$, and $\frac1{m^2}\le\frac2{m(m+1)}$. If $\Delta\le\frac1{16}$,
take $k=\lceil m\delta\rceil$, so that $k\le\frac m4+1\le\frac m2$. Then
$\frac m{m-k}\le1+\frac{2k}m\le1+2\delta+\frac2m\le1+4\delta$, because $\frac1m\le\Delta\le\delta$.
Also $\frac\Delta{k+1}\le\frac\delta m$. So
$f_N(m)\le X(1+4\delta)+\frac{3242x\delta}{m^2}\le X(1+4\delta+6484\delta)$, which gives the
claim. If $\Delta>\frac1{16}$, take $k=\lfloor m/2\rfloor$, so that $m-k\ge\frac m2$ and
$k+1\ge\frac m2$. Then $f_N(m)\le2X+\frac{6484x\Delta}{m^2}\le X(2+12968\Delta)$. Here
$\frac14<\delta\le1.073$, so $1<4\delta$ and $\Delta\le1.073\,\delta$. Hence
$2+12968\Delta\le1+4\delta+13915\delta$, which again gives the claim.
\end{proof}

\begin{proof}[Proof of Theorem~\ref{thm:profile}]
We check that Propositions~\ref{thm:B-local-lower}
and~\ref{thm:B-local-upper} apply on the whole range of $m$ and that their
error terms tend to $0$ uniformly there. Then we show that the mirror term
$f_N(N-m)$ is negligible.

\emph{Step 1. The hypotheses hold.} Put $\ell_x=\log(x+1)\ge\log2$ and let
$\mathcal R_N$ be the range
$\omega_Nx\ell_x^2\le m\le\delta'_NN$. Since $x\log^2N/N\to0$, for large $N$
we have $\varepsilon(1+\log N)\le\min(\bar c,\frac1{100})$,
$\omega_N\ge\max(A_1,16)$, $\delta'_N\le\frac18$ and $804x\log(x+2)\le N$.
So \eqref{eq:H} holds as well. If $\mathcal R_N$ is empty there is nothing to prove.

\emph{Step 2. The errors are small on $\mathcal R_N$.} The function
$\frac{x\log^2m}m$ decreases for $m\ge e^2$. At
$m=\omega_Nx\ell_x^2$ we have $m\ge e^2$ and $\log m\le\log\omega_N+3\ell_x$,
as $\log x$ and $\log\ell_x$ are at most $\ell_x$. So on $\mathcal R_N$, uniformly in
$x\ge1$,
$\frac{x\log^2m}m\le\frac{(\log\omega_N+3\ell_x)^2}{\omega_N\ell_x^2}
\le\frac1{\omega_N}(\frac{\log\omega_N}{\log2}+3)^2\to0$. Also
$m\ge\omega_N\log^22\to\infty$ on $\mathcal R_N$. Since $\log(m+1)\le2\log^2m$ for
$m\ge3$, Proposition~\ref{thm:B-local-lower} has $e_1(m)\le
C(\frac{x\log^2m}m+\varepsilon\log N+me^{-m/9})\to0$ uniformly on $\mathcal R_N$. Also
$\Delta(m)\le\frac{x\log^2m}m+\varepsilon\log N+\delta'_N\to0$ there, and
Proposition~\ref{thm:B-local-upper} applies. So
$f_N(m)=\frac{x}{m(m+1)}(1+o(1))$ uniformly on $\mathcal R_N$.

\emph{Step 3. The mirror term.} Lemma~\ref{lem:B-mode} gives $m_h\le A_0x\ell_x^2\le h$
for large $N$, since $x\ell_x^2\le x\log^2(N+1)=o(N)$. As $N-m\ge h$,
Lemma~\ref{lem:unimodal} gives $0\le f_N(N-m)\le f_N(h)=C_N$. By
Theorem~\ref{thm:erw-bounds}(b) and $\varepsilon\log N\to0$,
$C_N\le\frac{4x}{N(N+2)}N^{2\varepsilon}e^{5200\varepsilon}=\frac{4x}{N^2}(1+o(1))$.
On $\mathcal R_N$ this is at most $\frac{x}{m(m+1)}\cdot8\delta_N'^2(1+o(1))$, since
$m(m+1)\le2\delta_N'^2N^2$. Hence
$P_N(N-m)=\frac12[f_N(m)+f_N(N-m)]=\frac{x}{2m(m+1)}(1+o(1))$ uniformly
on $\mathcal R_N$, and $m(m+1)=m^2(1+o(1))$ because $m\to\infty$.
\end{proof}

\section{Odd \texorpdfstring{$N$}{N}}\label{or:odd}

We show that the crossing for odd $N$ behaves as for even $N$. As in
Proposition~\ref{prop:odd}, $N\ge3$ is odd, $c=1-2\varepsilon$, $\widetilde C_N$ is the
probability of each of the two central values, and $x_N=N\varepsilon_N$ is the crossing.

\begin{proposition}[odd $N$, the crossing]\label{prop:odd-more}
\begin{enumerate}[(a)]
\item As $N\to\infty$ over odd $N$, $x_N+\log x_N=\Lambda_N+O(L^2/N)$. So both parts of
Theorem~\ref{thm:erw-crossing} hold for odd $N$ as well. More precisely, if $L\ge30$, then
$\lvert x_N-g_N\rvert\le\tilde\rho_N$ and $\lvert x_N+\log x_N-\Lambda_N\rvert\le\tilde\rho_N$,
where $\tilde\rho_N=(24L^2+20812L)/N\le1/L^2$. In particular the last claim of
Corollary~\ref{cor:mode-location}(b) holds at the crossing for every large odd $N$.
\item For every odd $N\ge1689$ we have $x_N>2$, and at the crossing $P_N(0)<P_N(1)$. So
Corollary~\ref{cor:one-point} applies to these $N$ with $x=x_N$.
\item There is $N_1'$ such that the law at the crossing is convex between the humps for
every odd $N\ge N_1'$. We do not compute the threshold $N_1'$.
\end{enumerate}
\end{proposition}

\begin{proof}
(a) We follow the proof of Theorem~\ref{thm:erw-crossing}(a), with
$\widetilde F(x)=\log\widetilde C_N-\log E_N$ at $\varepsilon=x/N$ in place of $F$. Let
$1\le x\le2\Lambda_N$. For large $N$, Theorem~\ref{thm:erw-bounds} applies to the even number
$N+1$ at this $\varepsilon$ and gives
$\log C_{N+1}=\log\frac{4\varepsilon}{N+3}+O(\varepsilon L)$. By Proposition~\ref{prop:odd}(a),
$\log\widetilde C_N=\log C_{N+1}+O(1/N)$. Here $\log\frac{4\varepsilon}{N+3}=\log x+\log4-2L+O(1/N)$,
and $1/N\le xL/N$. So
\[
 \log\widetilde C_N=\log x+\log4-2L+O\Bigl(\frac{xL}N\Bigr),
\]
which is the estimate for $\log C_N$ in Step~1 of that proof. The formula
$E_N=\frac12(1-\varepsilon)^{N-1}$ holds for every $N$ (Lemma~\ref{lem:erw-reduction}). So
\eqref{eq:F-G} holds for $\widetilde F$, and Step~2 of that proof gives
$\lvert x_N-g_N\rvert=O(L^2/N)$ and the claim. Part (b) of Theorem~\ref{thm:erw-crossing} follows
from part (a) as before.

For the explicit bound let $L\ge30$. As in the proof of Theorem~\ref{thm:erw-crossing},
$e^L\ge L^{12}/12!$ gives $26L^4+20802L^3\le N$. Let $1\le x\le4L$ and $\varepsilon=x/N$. Since
$\log(N+1)\le L+1$, we get $\varepsilon(1+\log(N+1))\le4L(L+2)/N\le\frac1{100}$, so
Theorem~\ref{thm:erw-bounds} applies to $N+1$. By Proposition~\ref{prop:odd}(a),
$\log\widetilde C_N=\log C_{N+1}-\log(1-u)$ with $u=c/N\le\frac1N$, and
$0\le-\log(1-u)\le u+u^2\le\frac2N$. Next
$\log\frac{4\varepsilon}{N+3}=\log x+\log4-2L-\log(1+\frac3N)$ with
$0\le\log(1+\frac3N)\le\frac3N$. Theorem~\ref{thm:erw-bounds} gives
$\log(1-2\varepsilon)\le\log C_{N+1}-\log\frac{4\varepsilon}{N+3}\le2\varepsilon\log(N+1)+5200\varepsilon$.
Here $\log(1-2\varepsilon)\ge-2\varepsilon-4\varepsilon^2$, $4\varepsilon^2\le\frac1N$,
$2\varepsilon\log(N+1)\le\frac{8L^2+8L}N$ and $5200\varepsilon\le\frac{20800L}N$. For the end,
$\log E_N=(N-1)\log(1-\varepsilon)-\log2$ lies between $-\log2-x-\frac{16L^2}N$ and
$-\log2-x+\varepsilon$, by $-y-y^2\le\log(1-y)\le-y$. Adding up, and using
$\varepsilon\le\frac{4L}N$, we find that $\widetilde F(x)-G(x)+\Lambda_N$ lies between
$-\frac{12L+4}N$ and $\frac{24L^2+20808L+2}N$. So \eqref{eq:F-G} holds for $\widetilde F$ with
$\tilde\rho_N=(24L^2+20812L)/N$. Then $\tilde\rho_N\le1/L^2$, both ranges in Step~2 of that
proof are nonempty, and Step~2 gives $\lvert x_N-g_N\rvert\le\tilde\rho_N$ and
$\lvert G(x_N)-\Lambda_N\rvert\le\tilde\rho_N$.

For the last claim, $x_N=2\log N-\log\log N+O(1)$ for odd $N$, so
$\log N\le x_N\le2\log N$ for every large odd $N$, and Corollary~\ref{cor:mode-location}(b)
applies with $x=x_N$.

(b) We argue as in the proof of Proposition~\ref{prop:dip-hyp}. Let $N\ge1689$ be odd and take
$\varepsilon=2/N$. Then $\varepsilon(1+\log(N+1))\le\frac1{100}$, so
Theorem~\ref{thm:erw-bounds}(b) for $N+1$ and Proposition~\ref{prop:odd}(a) give
\[
 \widetilde C_N\le\frac N{N-1}\,C_{N+1}\le\frac N{N-1}\cdot\frac8{N(N+3)}(N+1)^{4/N}e^{10400/N}.
\]
The right side decreases in $N$ and is below $1.4\cdot10^{-3}$ at $N=1689$. Also
$E_N=\frac12(1-\frac2N)^{N-1}>0.06$, as in that proof. So $\widetilde C_N<E_N$ at
$\varepsilon=2/N$, and hence $\varepsilon_N>2/N$ and $x_N>2\ge2(1-\varepsilon_N)$.
Theorem~\ref{thm:dip} then gives $P_N(N-1)>E_N$ at the crossing, which is $P_N(0)<P_N(1)$ by
symmetry. This is the hypothesis of the second case of Corollary~\ref{cor:one-point}.

(c) Proposition~\ref{thm:convex-large} holds for every large $N$, even or odd, and every
$\varepsilon$ with $\log N\le x\le2\log N$. By (a), $\log N\le x_N\le2\log N$ for every large
odd $N$. So Proposition~\ref{thm:convex-large} applies at the crossing.
\end{proof}

\section{A likelihood for reversible fluctuation assays}\label{or:likelihood}

This section proves likelihood bounds for the reversible fluctuation
assay of Bokes, Hlubinov\'a and Singh~\cite{bhs2023}. It uses the notation of
Section~\ref{sec:results} of the paper and the bound $\delta_\Sigma$ of
Theorem~\ref{thm:cp}.

With one sensitive founder, uniform parent choice, and independent
daughter-type switching with probability $\varepsilon$ at each birth in
either direction, the resistant count at population $t$ increases from
$k$ with probability $\varepsilon+(1-2\varepsilon)k/t$. Its law at size
$N$ is therefore $f_N$. Write $P_{N,x}$ for this law at
$\varepsilon=x/N$, and $Q_{N,x}=\CP_N$ with masses $q_{N,x}(k)$.
The CP law truncates jump sizes, not the total count. Its likelihood
uses the full law on $\mathbb N_0$.

\begin{theorem}[constrained compound Poisson likelihood]\label{thm:assay-mle}
Let $N\ge2,m\ge1$ be integers, $a_N=1-1/N$, and fix in advance
$\Theta=[\ell,u]$ with $0<\ell<u<N$ and $a_Nu>1$.
Let $K_1,\ldots,K_m$ be independent observations with law $P_{N,x_0}$,
where $x_0\in\Theta$. Define $\widehat x$ as the smallest maximizer of
$L_x=\prod_{i=1}^m q_{N,x}(K_i)$ on $\Theta$, and put
\[
 \begin{gathered}
 W=u-\ell,\qquad \omega=(ma_N)^{-1},\qquad
 J=\lceil W/\omega\rceil+1,\\
 v_x=\min\{1,\delta_\Sigma(x)\},\qquad
 \eta_x=1-(1-v_x)^m\le\min\{1,mv_x\},
 \end{gathered}
\]
where $\delta_\Sigma(x)$ is evaluated at $\varepsilon=x/N$.
Probabilities and expectations below are under $P_{N,x_0}^{\otimes m}$.
For every $r>\omega$,
\begin{equation}\label{eq:assay-tail}
 \Prob\{|\widehat x-x_0|\ge r\}
 \le\min\left\{1,\eta_{x_0}+e^{1/2}J
       \exp\left(-\frac{m(r-\omega)^2}{2e^2u^2}\right)\right\}.
\end{equation}
For $0<\alpha<1$, set
$r_\alpha=\omega+eu\sqrt{2\log(e^{1/2}J/\alpha)/m}$. Then
\begin{equation}\label{eq:assay-bias}
 |\E\widehat x-x_0|\le\E|\widehat x-x_0|
 \le\min\{W,r_\alpha+W(\alpha+\eta_{x_0})\}.
\end{equation}
In particular, for fixed $0<a<b$, $\Theta=[a\log N,b\log N]$ and
$m=\lceil(\log N)^3\rceil$, uniformly in $x_0\in\Theta$ as $N\to\infty$,
\[
 \E|\widehat x-x_0|
 =O_{a,b}\left(\sqrt{\frac{\log\log N}{\log N}}
              +\frac1{\log N}+\frac{\log^8N}{N}\right)=o(1).
\]
Thus $\widehat x$ is absolutely consistent in probability and its mean
bias tends to zero uniformly. For $\widehat\varepsilon=\widehat x/N$,
divide \eqref{eq:assay-bias} by $N$.
\end{theorem}

The estimator is constrained and uses a sensitive founder. Its mean-bias
bound is absolute, without a signed bias expansion.

\begin{proof}
For each integer $k\ge0$, $q_{N,x}(k)$ is $e^{-a_Nx}$ times a finite
polynomial in $x$, and is positive, since the size-one jumps alone contribute
$e^{-a_Nx}(x/2)^k/k!$. Thus the likelihood attains its maximum on
$\Theta$. Its smallest maximizer is measurable on the countable sample
space $\mathbb N_0^m$.

Put $G_N(s)=\sum_{d=1}^{N-1}(1-s^d)/[d(d+1)]$.
This function decreases continuously and strictly from $a_N$ to $0$ on
$[0,1]$. Since $a_Nu>1$, choose $s\in(0,1)$ with $G_N(s)=1/u$.
The statistic $s^K\in[0,1]$ has mean $e^{-x/u}$ under $Q_{N,x}$, hence
$\operatorname{TV}(Q_{N,x},Q_{N,y})\ge e^{-1}|x-y|/u$ for $x,y\in\Theta$.
For the affinity $\rho(x,y)=\sum_k\sqrt{q_{N,x}(k)q_{N,y}(k)}$,
Cauchy--Schwarz gives $\operatorname{TV}^2\le1-\rho^2$, and therefore
\[
 \rho(x,y)\le\exp\{-|x-y|^2/(2e^2u^2)\}.
\]
The CP semigroup gives
$q_{N,y}(k)\ge e^{-a_N(y-x)}q_{N,x}(k)$ for $y\ge x$.
Choose $J$ equally spaced grid points including $\ell,u$, with spacing at
most $\omega$. If $|\widehat x-x_0|\ge r$, the grid point $y$ at or
immediately above $\widehat x$ has $|y-x_0|\ge r-\omega$ and
$L_y\ge e^{-ma_N\omega}L_{\widehat x}\ge e^{-1}L_{x_0}$.
For each fixed grid point $y$, under $Q_{N,x_0}^{\otimes m}$, Markov's inequality gives
\[
 \Prob\{L_y/L_{x_0}\ge e^{-1}\}
 \le e^{1/2}\E\sqrt{L_y/L_{x_0}}
 =e^{1/2}\rho(x_0,y)^m.
\]
The union bound over the grid proves \eqref{eq:assay-tail} without
$\eta_{x_0}$ for CP observations. Independent couplings from
Theorem~\ref{thm:cp} give
$\operatorname{TV}(P_{N,x_0}^{\otimes m},Q_{N,x_0}^{\otimes m})
\le1-(1-v_{x_0})^m=\eta_{x_0}$, proving the stated bound.
Substitute $r=r_\alpha$ and use $|\widehat x-x_0|\le W$ to split its
expectation at $r_\alpha$, which proves \eqref{eq:assay-bias}.

Finally, with $L=\log N$, the stated interval and sample size give
$\sup_{x\in\Theta}\delta_\Sigma(x)=O_{a,b}(L^4/N)$,
$\omega=O(L^{-3})$, $J=O_{a,b}(L^4)$ and
$\sup_{x\in\Theta}\eta_x=O_{a,b}(L^7/N)$.
Taking $\alpha=L^{-2}$ in \eqref{eq:assay-bias} proves the uniform
mean-error bound. Markov's inequality gives absolute consistency.
\end{proof}

\section{Convexity between the humps}\label{or:window}

This section concerns a stronger, prescribed convexity window between the
two humps of the law at the crossing. We prove it for every large $N$,
with a threshold $N_1$ that we do not compute. A certified computation
gives it at the crossing for $5\le N\le2000$, and it is open in between
(Conjecture~\ref{conj:convex}). Proposition~\ref{prop:main-window} collects
these statements. The shape of Theorem~\ref{thm:main-shape}(b) does not
need this window. For a mode $a^*\ge N/2$ of $P_N$ we put $m^*=N-a^*$, as
in Section~\ref{sec:shape}.

\begin{proposition}[an eventual convexity window]\label{prop:main-window}
There is $N_1$ such that the following holds for every $N\ge N_1$, even
or odd, and every $\varepsilon$ with $\log N\le x\le2\log N$. For every
mode $a^*\ge N/2$ of $P_N$ we have $D_N(a)>0$ for
$2m^*+1\le a\le N-2m^*-1$ (Proposition~\ref{thm:convex-large}). In
particular the window holds at the crossing for every large even $N$.
We do not compute $N_1$. At the crossing, a certified computation gives
the window for every $5\le N\le2000$
(Proposition~\ref{prop:convex-certified}). For $N>2000$ it is a conjecture
(Conjecture~\ref{conj:convex}).
\end{proposition}

\begin{definition}[convex between the humps]\label{def:convex}
Let $N\ge5$. The law at the crossing is \emph{convex between the humps} if,
for $x=x_N$ and every mode $a^*\ge N/2$ of $P_N$ with $m^*=N-a^*$, the second
difference $P_N(a-1)-2P_N(a)+P_N(a+1)$ is positive for
$2m^*+1\le a\le N-2m^*-1$.
\end{definition}

With the symmetry of the law, this makes $P_N$ strictly convex on
$[2m^*,N-2m^*]$, so $h$ is its unique minimizer there (the two central values
for odd $N$).

\begin{corollary}[one inequality is enough]\label{cor:one-point}
Let $N\ge8$ be even and $x=x_N$, or let $N\ge5$ and $0<\varepsilon<\frac12$
be such that $P_N(0)<P_N(1)$. Let $a^*$ be any mode of $P_N$ and
$m^*=N-a^*$. If $D_N(a_1)>0$ for some integer $a_1$ with $m^*<a_1\le2m^*+1$
and $a_1\le N/2$, then $D_N(a)>0$ for every $2m^*+1\le a\le N-2m^*-1$. In
particular $P_N$ is strictly convex on $[2m^*,N-2m^*]$ as soon as
$D_N(2m^*+1)>0$.
\end{corollary}

\begin{proof}
At the crossing with $N\ge8$ even, $P_N(0)=E_N<P_N(N-1)=P_N(1)$ by
Proposition~\ref{prop:dip-hyp}. Let $a_0$ be the smallest maximizer of $P_N$.
Then $1\le a_0\le m^*$, because $P_N(0)<P_N(1)$ and $P_N(m^*)=P_N(a^*)$ is the
maximum. Also $P_N(a_0-1)<P_N(a_0)\ge P_N(a_0+1)$, so $D_N(a_0)<0$. Now
Theorem~\ref{thm:no-holes} applies to $a_0<a_1$.
\end{proof}

The condition $a_1\le N/2$ costs nothing at $a_1=2m^*+1$, because
$2m^*+1\le N/2$ whenever the range $2m^*+1\le a\le N-2m^*-1$ is not empty.
If $a^*<N/2$, then $m^*>N/2$, and no $a_1$ meets the conditions.

\begin{proposition}[convexity between the humps for large $N$]\label{thm:convex-large}
There is $N_1$ such that the following holds for every $N\ge N_1$, even or odd, and every
$\varepsilon$ with $\log N\le x\le2\log N$. For every mode $a^*\ge N/2$ of $P_N$ with
$m^*=N-a^*$ we have $4m^*+4\le N$ and
\[
 P_N(a-1)-2P_N(a)+P_N(a+1)>0\qquad\text{for }2m^*+1\le a\le N-2m^*-1 .
\]
In particular, after enlarging $N_1$, the law at the crossing is convex between the humps for
every even $N\ge N_1$. We do not compute the threshold $N_1$ (Remark~\ref{rem:convex-open}).
\end{proposition}

\paragraph{Idea of the proof.}
If there were no point of strict convexity between $m^*$ and
$2m^*+1$, concavity would force more than half the total mass
into an interval strictly to the left of the center.
Symmetry rules this out, and Corollary~\ref{cor:one-point}
then gives the prescribed window.

\begin{proof}
Let $\log N\le x\le2\log N$ and let $a^*\ge N/2$ be any mode.
Write $m=m^*=N-a^*$ and $M=P_N(m)$. All estimates below are
uniform in $x$ and in the choice of mode, for either parity of $N$.
By Corollary~\ref{cor:mode-location}(b),
\[
 m=x\log x+\lambda_{\mathrm L}x+O(\sqrt{x\log x}),
 \qquad \lambda_m\longrightarrow\lambda_{\mathrm L}.
\]
The error in Theorem~\ref{thm:landau}(b) tends uniformly to zero,
since $(1+\log x)/x=o(1)$ and $x\delta_B=O(\log^5N/N)$.
Consequently
\[
 2xM\ge xf_N(m)=f_L(\lambda_m)+o(1)
                =f_L(\lambda_{\mathrm L})+o(1).
\]
It follows that $M(m+3)>1$, $m\ge1$ and $2m+2\le N/2$
for all sufficiently large $N$.

Suppose that $D_N(a)\le0$ for every integer $m<a\le2m+1$.
Then $P_N$ is concave on $[m,2m+2]$, and $P_N(2m+2)\ge0$ gives
\[
 \sum_{a=m}^{2m+1}P_N(a)
 \ge\frac{M}{m+2}\sum_{j=1}^{m+2}j
 =\frac{M(m+3)}2>\frac12.
\]
But this interval and its reflection are disjoint and have equal
mass, a contradiction. Thus $D_N(a_1)>0$ for some
$m<a_1\le2m+1\le N/2$.
Theorem~\ref{thm:dip} gives $P_N(0)<P_N(1)$ because
$x\ge\log N>2$, and $\varepsilon<1/2$ for large $N$.
Corollary~\ref{cor:one-point} now proves the required positivity.
The final assertion follows from Theorem~\ref{thm:erw-crossing}(b).
\end{proof}

\begin{proposition}[small $N$, computer-assisted]\label{prop:convex-certified}
For every $5\le N\le N_{\rm c}=2000$, even or odd, the law at the crossing is
convex between the humps.
\end{proposition}

\begin{proof}
For $10\le N\le2000$, the certified computation described in
Section~\ref{sec:verify} brackets $\varepsilon_N$ and encloses $P_N$ throughout
that bracket. It lists every possible mode and checks $D_N(a)>0$ on each
required range. When the range is nonempty, it also checks
$D_N(a_0)<0$ at some $1\le a_0\le2m^*$ and $D_N(2m^*+1)>0$.
Theorem~\ref{thm:no-holes} then gives a second check of the conclusion.

For $5\le N\le9$, Lean brackets the crossing by exact rational
arithmetic and Proposition~\ref{prop:mlr}. Over the whole bracket it checks
$\varepsilon<\frac12$ and $P_N(0)<P_N(1)$. For each possible mode it
checks either that the prescribed window is empty, or that
$2m^*+1\le N/2$ and $D_N(2m^*+1)>0$.
Corollary~\ref{cor:one-point} then gives the claim. These are
outward-rounded integer computations with kernel-checked soundness
(Section~\ref{sec:verify}). At $N=5$ the only possible right-hand mode is
$a^*=4$, whose window is empty.
\end{proof}

In double precision the law at the crossing is also convex between the humps
for $N=3200$, $6400$, $12800$ and $25600$. The definition leaves out the two
endpoints $a=2m^*$ and
$a=N-2m^*$, because there the second difference is not positive for some
$N\le51$. Heuristically, between the humps $P_N(a)$ is close to
$\frac x2\bigl[\frac1{a(a+1)}+\frac1{(N-a)(N-a+1)}\bigr]$, which is strictly
convex in $a$. We also leave open the tail for $m$ between about $x\log x$
and $x\log^2x$, below the range of the tail theorem in
\arxivonly{Theorem}{thm:profile}.

The constants of the large-$N$ result above are not computed, so the
following is open.

\begin{conjecture}[convexity between the humps]\label{conj:convex}
For every $N>2000$ the law at the crossing is convex between the humps.
\end{conjecture}

\begin{remark}[what is still open]\label{rem:convex-open}
The remaining gap concerns the prescribed convexity window.
The minimum and two-hump shape hold for every even $N\ge8$,
and Corollary~\ref{cor:height-window} gives a nonempty convexity
window there. The threshold $N_1$ in Proposition~\ref{thm:convex-large}
is not computed, and its proof uses mode location and the local
mass bound to obtain $P_N(m^*)(m^*+3)>1$. At the crossing
$m^*\approx x\log x$ and $P_N(m^*)\approx f_L(\lambda_{\mathrm L})/(2x)$, so
this condition needs roughly $f_L(\lambda_{\mathrm L})\log x>2$, that is $x$
larger than about $6.4\cdot10^4$. Since $x_N\approx2\log N$, no computation
can reach this. In double precision the product $P_N(m^*)(m^*+3)$ is
about $0.235$, $0.239$ and $0.24$ at $N=1000$, $2000$ and $4000$, so it
grows very slowly. These numbers are heuristic and are not used in any proof.
\end{remark}

\section{Proofs for the threshold certificate}\label{or:threshold}

This section proves the lemmas that Appendix~\ref{app:threshold} summarizes,
and so completes the proof of Theorem~\ref{thm:threshold}. We use the notation
of that appendix, and in particular the inequality~\eqref{eq:thr-fixed}. We
write $\Phi_N(t)=\sum_{n\le N}c_n(it)^n$ for the partial sums of the moment
series.

\subsection{One value of the memory}

\begin{lemma}[moment bounds]\label{lem:thr-moments}
Let $N\ge3$.
\begin{enumerate}[(a)]
\item Put $A=((N+1)a-1)/N$. If $m_j\le AB^j$ for $1\le j\le N$, then
$m_n\le AB^n$ for all $n\ge1$.
\item Let $S\subset\mathbb C$ be compact, with $\frac12<\alpha\le\operatorname{Re}z$,
$\lvert\operatorname{Im}z\rvert\le\rho$ and $\lvert z\rvert\le1$ for $z\in S$.
Define $m_n(z)$ by the recursion with $z$ in place of $a$, and
$c_n(z)=m_n(z)/\Gamma(1+nz)$. Then every $m_n$ is analytic near $S$. Put
$A=((N+1)\alpha-1)/N$. If $\lvert m_j\rvert\le AB^j$ on $S$ for $1\le j\le N$,
then for all $n\ge1$ and $z\in S$
\[
 \lvert m_n(z)\rvert\le AB^n,\qquad
 \lvert c_n(z)\rvert\le\frac{\tilde A\tilde B^n}{\Gamma(1+n\alpha)},\qquad
 \tilde A=Ae^{\rho^2/(2\alpha^2)},\quad \tilde B=Be^{\rho^2/(2\alpha)} .
\]
\end{enumerate}
\end{lemma}

\begin{proof}
For (a) we use induction on $n>N$. As $\theta_j\le1$ and
$na-\theta_n\ge na-1>0$, the bound for smaller indices gives
$m_n\le(n-1)A^2B^n/(na-1)$. The ratio $(na-1)/(n-1)$ increases in $n$, since
its derivative is $(1-a)/(n-1)^2$, and it equals $A$ at $n=N+1$. So
$m_n\le AB^n$.

For (b), $\operatorname{Re}(nz-\theta_n)\ge n\alpha-1>0$ on $S$, because
$\alpha\le1$. So $m_n$ is analytic near $S$, and the induction of (a) goes
through with $\lvert\theta_j\rvert\le1$. For $x>0$ and real $y$,
\[
 \Bigl\lvert\frac{\Gamma(x)}{\Gamma(x+iy)}\Bigr\rvert^2
 =\prod_{k\ge0}\Bigl(1+\frac{y^2}{(x+k)^2}\Bigr)
 \le\exp\Bigl(\frac{y^2}{x^2}+\frac{y^2}x\Bigr)
\]
by~\cite[5.8.3]{dlmf}. We take $x=1+n\operatorname{Re}z\ge1+n\alpha$ and
$\lvert y\rvert\le n\rho$. Then $y^2/x^2\le\rho^2/\alpha^2$ and
$y^2/x\le n\rho^2/\alpha$. Finally
$\Gamma(1+n\operatorname{Re}z)\ge\Gamma(1+n\alpha)$, as $\Gamma$ increases on
$[1.47,\infty)$.
\end{proof}

\begin{lemma}[truncation]\label{lem:thr-trunc}
Let $0<\alpha\le1$, and suppose that
$\lvert c_n\rvert\le\tilde A\tilde B^n/\Gamma(1+n\alpha)$ for $n>N$. Let $t>0$,
$0\le k\le N$ and $n_0=N+1$, and put
\[
 u=\frac{\tilde A\tilde B^{n_0}n_0^kt^{n_0-k}}{\Gamma(1+\alpha n_0)},\qquad
 r=\tilde Bt\Bigl(1+\frac1{n_0}\Bigr)^k
 \frac{\Gamma(1+\alpha n_0)}{\Gamma(1+\alpha n_0+\alpha)} .
\]
If $r<1$, then
$\bigl\lvert\frac{d^k}{dt^k}\sum_{n>N}c_n(it)^n\bigr\rvert\le\tau_k(t):=u/(1-r)$,
and $\tau_k$ increases on $(0,t]$.
\end{lemma}

\begin{proof}
The $k$th derivative of $(it)^n$ has modulus at most $n^kt^{n-k}$. The ratio
of consecutive terms of the majorant
$\tilde A\tilde B^nn^kt^{n-k}/\Gamma(1+\alpha n)$ is
$\tilde Bt(1+1/n)^k\Gamma(1+\alpha n)/\Gamma(1+\alpha n+\alpha)$. It decreases
in $n$, because $\log\Gamma$ is convex. So the majorant is at most a
geometric series with first term $u$ and ratio $r$. Both $u$ and $r$ increase
in $t$.
\end{proof}

\begin{lemma}[cells]\label{lem:thr-cells}
Let $\nu$ be even, $t_c>0$ and $\lvert\delta\rvert\le h$. Then
$\lvert\phi(t_c+\delta)\rvert\le\sum_{k<\nu}\lvert\phi^{(k)}(t_c)\rvert
h^k/k!+c_\nu h^\nu$.
\end{lemma}

\begin{proof}
This is Taylor's formula, with
$\lvert\phi^{(\nu)}\rvert\le\E\lvert L\rvert^\nu=\E L^\nu=\nu!\,c_\nu$.
\end{proof}

\begin{lemma}[the tail]\label{lem:thr-tail}
Let $0<T_2<T_3$ with $T_3\ge2^aT_2$, and let $\beta>3$ with $a\beta>1$.
Cover $[0,T_3]$ by cells $[l_j,r_j]$ with $l_0=0$ and with $T_2$ among the
end points, and let $s_j\ge\sup_{[l_j,r_j]}\lvert\phi\rvert$. Let $C$ be the
largest $s_jr_j^\beta$ over the cells that meet $[T_2,T_3]$, and put
\[
 R=T_2^{-1/a}\sum_{r_j\le T_2}s_j
 \Bigl(1-\frac12\Bigl(\frac{r_j}{T_2}\Bigr)^{1/a}\Bigr)^{-a\beta}
 \bigl(r_j^{1/a}-l_j^{1/a}\bigr)+\frac{2^{a\beta}CT_2^{-\beta}}{a\beta-1} .
\]
If $R<1$, then $\lvert\phi(t)\rvert\le Ct^{-\beta}$ for every $t\ge T_3$.
\end{lemma}

\begin{proof}
Let $b=s_j$ on the cells in $[0,T_3)$ and $b(t)=Ct^{-\beta}$ for $t\ge T_3$.
On a cell that meets $[T_2,T_3]$ we have $s_j\le Cr_j^{-\beta}\le Ct^{-\beta}$.
So $b(t)\le Ct^{-\beta}$ for $t\ge T_2$, and $\lvert\phi\rvert\le b$ on
$[0,T_3]$. Let $J(t)=\int_0^1b(tv^a)\,b(t(1-v)^a)\,dv$.

We show that $J(t)\le Rb(t)$ for $t\ge T_3$. By symmetry $J(t)$ is twice the
integral over $v\le\frac12$. There $t(1-v)^a\ge2^{-a}T_3\ge T_2$, so
$b(t(1-v)^a)\le b(t)(1-v)^{-a\beta}$. Hence $J(t)/b(t)$ is at most
$2\int_0^{1/2}b(tv^a)(1-v)^{-a\beta}\,dv$. Put $v_0=(T_2/t)^{1/a}\le\frac12$.
On $v<v_0$ we substitute $u=tv^a$, and this part becomes
\[
 \frac2at^{-1/a}\int_0^{T_2}s(u)\bigl(1-(u/t)^{1/a}\bigr)^{-a\beta}
 u^{1/a-1}\,du .
\]
Here $t^{-1/a}\le\frac12T_2^{-1/a}$ and $(u/t)^{1/a}\le\frac12(u/T_2)^{1/a}$.
The weight increases in $u$, so on each cell we take it at $r_j$. As
$\int_l^ru^{1/a-1}du=a(r^{1/a}-l^{1/a})$, this part is at most the first term
of $R$. On $v\ge v_0$ we use $b(tv^a)\le Ct^{-\beta}v^{-a\beta}$ and
$(1-v)^{-a\beta}\le2^{a\beta}$. This part is at most
$2^{1+a\beta}T_2^{-\beta}(T_2/t)^{1/a}C/(a\beta-1)$, which is at most the
second term of $R$.

Let $\varsigma$ be the supremum of the $\tau\ge T_3$ with
$\lvert\phi\rvert\le b$ on $[0,\tau]$, and suppose that $\varsigma<\infty$. By
continuity $\lvert\phi\rvert\le b$ on $[0,\varsigma]$. The arguments in
\eqref{eq:thr-fixed} at $t=\varsigma$ lie in $[0,\varsigma]$, so
$\lvert\phi(\varsigma)\rvert\le J(\varsigma)\le Rb(\varsigma)<b(\varsigma)$.
As $\lvert\phi\rvert$ and $b$ are continuous on $[T_3,\infty)$, we get
$\lvert\phi\rvert<b$ near $\varsigma$, against the choice of $\varsigma$.
\end{proof}

This is an explicit form of Step~4 in the proof of~\cite[Theorem~2.5]{glr2026},
with computed constants.

\begin{lemma}[the second derivative]\label{lem:thr-fpp}
If $\lvert\phi(t)\rvert\le Ct^{-\beta}$ for $t\ge T_3$, with $\beta>3$, then
$g_a$ has two continuous derivatives and
\[
 g_a''(0)=-\frac1\pi\Bigl(\sum_{\substack{n\le N\\ n\text{ even}}}
 (-1)^{n/2}c_n\frac{T_3^{n+3}}{n+3}+E\Bigr),\qquad
 \lvert E\rvert\le\frac{T_3^3}3\tau_0(T_3)+\frac{CT_3^{3-\beta}}{\beta-3} .
\]
\end{lemma}

\begin{proof}
As $t^2\operatorname{Re}\phi(t)$ is integrable, Fourier inversion gives
$g_a(x)=\frac1{2\pi}\int e^{-itx}\operatorname{Re}\phi(t)\,dt$, and dominated
convergence allows two derivatives under the
integral~\cite[Section~3.3]{durrett2019}. So
$g_a''(0)=-\frac1\pi\int_0^\infty t^2\operatorname{Re}\phi(t)\,dt$. On
$[0,T_3]$ we replace $\phi$ by $\Phi_N$, whose real part is
$\sum_{n\text{ even}}(-1)^{n/2}c_nt^n$, and integrate term by term. By
Lemma~\ref{lem:thr-trunc} the error is at most
$\int_0^{T_3}t^2\tau_0(t)\,dt\le\frac13T_3^3\tau_0(T_3)$. The integral over
$t>T_3$ is at most $CT_3^{3-\beta}/(\beta-3)$.
\end{proof}

\subsection{Intervals of the memory}

Ball arithmetic with $a$ as an interval fails here. Near the threshold the
series for $\phi$ cancels terms of size about $10^{60}$, so the boxes in $a$
would have to be narrower than about $10^{-65}$. We expand in $a$ around one
exact point instead.

\begin{lemma}[Cauchy remainder]\label{lem:thr-taylor}
Let $h$ be analytic near the closed disc $\lvert z-a_0\rvert\le\rho$, with
$\lvert h\rvert\le M$ there and Taylor coefficients $h_i$ at $a_0$. If $0<q<1$
and $\lvert a-a_0\rvert\le q\rho$, then
$\lvert h(a)-\sum_{i\le K}h_i(a-a_0)^i\rvert\le Mq^{K+1}/(1-q)$.
\end{lemma}

\begin{proof}
Cauchy's estimate gives $\lvert h_i\rvert\le M\rho^{-i}$. Sum the geometric
tail.
\end{proof}

Let $I=[a_0-w,a_0+w]$, $0<q<1$ and $\rho=w/q$, and let $S$ be the square
$a_0+[-\rho,\rho]+i[-\rho,\rho]$, which contains the disc of
Lemma~\ref{lem:thr-taylor}. Suppose that $\alpha=a_0-\rho>\frac12$, that
$\lvert z\rvert\le1$ on $S$, and that $\tilde A,\tilde B$ are as in
Lemma~\ref{lem:thr-moments}(b). Let $\tau_k$ be the bound of
Lemma~\ref{lem:thr-trunc} with these constants. It holds for every real
$a\in I$. For a cell with centre $t_c$ write
$\partial_t^k\Phi_N(z,t_c)=E_k(z)+iO_k(z)$, where $E_k$ is the sum over even
$n$ and $iO_k$ the sum over odd $n$. Both $E_k$ and $O_k$ are real for real
$z$. Let $E_{k,i}$ and $O_{k,i}$ be their Taylor coefficients at $a_0$, and
put $M_k=\sum_{n\le N}\sup_S\lvert c_n\rvert\,n(n-1)\cdots(n-k+1)\,t_c^{n-k}$.

\begin{proposition}[enclosure on an interval]\label{prop:thr-interval}
\begin{enumerate}[(a)]
\item For every $a\in I$ and every $t$ in the cell with centre $t_c$ and
half-width $h$,
\[
 \lvert\phi_a(t)\rvert\le\sum_{k<12}\Bigl(\sum_{i\le K}w^i
 \sqrt{E_{k,i}^2+O_{k,i}^2}+\frac{M_kq^{K+1}}{1-q}+\tau_k(t_c)\Bigr)
 \frac{h^k}{k!}+h^{12}\sup_S\lvert c_{12}\rvert .
\]
\item Let $T_2$, $T_3\ge2^{a_0+w}T_2$, $\beta$ and the cells be as in
Lemma~\ref{lem:thr-tail}, with $s_j$ the bound of (a), and suppose that $R<1$
for every $a\in I$. Let $F(z)=\sum_{n\le N,\,n\text{ even}}(-1)^{n/2}c_n(z)
T_3^{n+3}/(n+3)$, with Taylor coefficients $F_i$ at $a_0$, and let $M_F$ be
the same sum with $\sup_S\lvert c_n\rvert$ in place of $(-1)^{n/2}c_n(z)$. Then
for every $a\in I$
\[
 \Bigl\lvert g_a''(0)+\frac1\pi\sum_{i\le K}F_i(a-a_0)^i\Bigr\rvert\le\eta
 =\frac1\pi\Bigl(\frac{M_Fq^{K+1}}{1-q}+\frac{T_3^3}3\tau_0(T_3)
 +\frac{CT_3^{3-\beta}}{\beta-3}\Bigr),
\]
and $a\mapsto g_a''(0)$ is continuous on $I$.
\end{enumerate}
\end{proposition}

\begin{proof}
For (a) we apply Lemma~\ref{lem:thr-cells} with $\nu=12$, and
Lemma~\ref{lem:thr-trunc} to replace $\phi_a^{(k)}$ by $\partial_t^k\Phi_N$.
For real $a$ we have
$\lvert\partial_t^k\Phi_N(a,t_c)\rvert=\lvert(E_k(a),O_k(a))\rvert$.
Lemma~\ref{lem:thr-taylor} for $E_k$ and $O_k$, with
$\sup\lvert E_k\rvert+\sup\lvert O_k\rvert\le M_k$ on $S$, and the triangle
inequality in $\mathbb R^2$ give the bracket. Also
$c_{12}(a)\le\sup_S\lvert c_{12}\rvert$.

For (b), $T_3\ge2^aT_2$ for every $a\in I$, so Lemma~\ref{lem:thr-tail} applies
at each $a$. Then Lemma~\ref{lem:thr-fpp} and Lemma~\ref{lem:thr-taylor} for
$F$ give the enclosure. By Lemma~\ref{lem:thr-moments}(b) the series for
$\phi_a(t)$ converges uniformly for $a\in I$ and $t$ in compact sets. So
$\phi_a(t)$ is continuous in $a$. Moreover $t^2\lvert\phi_a(t)\rvert$ is at
most $t^2$ on $[0,T_3]$ and $Ct^{2-\beta}$ beyond, uniformly in $a\in I$.
Dominated convergence gives the continuity.
\end{proof}

Near $a=\frac12$ the bounds of Lemma~\ref{lem:thr-moments} degenerate, since
$m_2=a/(2a-1)$. We rescale.

\begin{lemma}[the scaled variable]\label{lem:thr-scaled}
Put $y=2a-1$ and $\mu_n=m_ny^{\lfloor n/2\rfloor}$.
\begin{enumerate}[(a)]
\item We have $\mu_0=\mu_1=1$, $\mu_2=a$, and for $n\ge3$
\[
 (na-\theta_n)\mu_n=\sum_{j=1}^{n-1}\theta_jy^{e_{nj}}\mu_j\mu_{n-j},
\]
where $e_{nj}=1$ if $n$ is even and $j$ is odd, and $e_{nj}=0$ otherwise.
Here $na-\theta_n$ is $(n-1)a$ for odd $n$ and $na-1$ for even $n$. So every
$\mu_n$ is a rational function of $a$ with no pole in
$\operatorname{Re}a>\frac14$.
\item Let $\sigma^2=\E L^2$, $X=L/\sigma$ and
$\hat c_n=\mu_n\Gamma(2a)^{\lfloor n/2\rfloor}/\Gamma(1+an)$. Then
$\sigma^2=1/(y\Gamma(2a))$,
\[
 \E e^{iuX}=\sum_{n\text{ even}}\hat c_n(iu)^n
 +\sqrt{y\Gamma(2a)}\sum_{n\text{ odd}}\hat c_n(iu)^n,
 \qquad
 g_a''(0)=-\frac1{\pi\sigma^3}\int_0^\infty u^2\operatorname{Re}\E e^{iuX}\,du .
\]
\end{enumerate}
\end{lemma}

\begin{proof}
Multiply the recursion for $m_n$ by $y^{\lfloor n/2\rfloor}$. The exponent
$\lfloor n/2\rfloor-\lfloor j/2\rfloor-\lfloor(n-j)/2\rfloor$ is $1$ when $n$ is
even and $j$ is odd, and $0$ otherwise. Also $m_2=a/y$. This proves (a), as
the zeros of the factors are $a=0$ and $a=1/n\le\frac14$. For (b),
$\sigma^2=2m_2/\Gamma(1+2a)=1/(y\Gamma(2a))$ and
$\E X^n/n!=c_n\sigma^{-n}=\mu_n\Gamma(2a)^{\lfloor n/2\rfloor}
(y\Gamma(2a))^{(n\bmod2)/2}/\Gamma(1+an)$. The formula for $g_a''(0)$ follows
from the substitution $t=u/\sigma$ in Lemma~\ref{lem:thr-fpp}.
\end{proof}

The fixed-point equation is linear, so $X$ satisfies it too, and
\eqref{eq:thr-fixed} and Lemmas~\ref{lem:thr-cells} to~\ref{lem:thr-fpp} hold
for $X$. The proof of Lemma~\ref{lem:thr-moments}(b) carries over to $\mu_n$
for $n\ge3$ and $\alpha>\frac14$, with two changes. The constant $A$ becomes
$((N+1)\alpha-1)/(N\kappa)$, where $\kappa$ is the largest value of
$\max(1,\lvert z\rvert)\max(1,\lvert2z-1\rvert)$ on $S$. And $\tilde B$ gains
the factor $\max(1,\sup_S\lvert\Gamma(2z)\rvert)^{1/2}$. For real $a$ the odd
part carries the factor $\sqrt{y\Gamma(2a)}$. In the cell bounds we treat the
even and odd parts as separate analytic functions. We multiply the odd part
by $f=\sup_I\sqrt{y\Gamma(2a)}$ only for real $a$, and the remainders by
$\max(1,f)$. The real part of $\E e^{iuX}$ involves only the even part. So
the disc may cross $a=\frac12$, and the cover reaches $a=\frac12$. The sign
of $g_a''(0)$ is the sign of the integral, since $\sigma^3>0$.

\subsection{From the runs to the theorem}\label{or:thr-runs}

The script accepts an interval of $a$ only when $G\pm\eta$ has one sign on
each piece of a subdivision of it. In the 7 scaled runs every piece of every
interval is negative, so $g_a''(0)<0$ on $(\frac12,0.668]$. The zero run has 8
intervals that cover $[0.668,0.6712]$. On the 2 intervals that cover
$[0.668,0.6688]$ every piece is negative. On the 5 intervals that cover
$[0.6692,0.6712]$ every piece is positive. On $[0.6688,0.6692]$ the pieces
up to $z_-$ are negative, the pieces from $z_+$ on are positive, and only
$(z_-,z_+)$ is undecided. Proposition~\ref{prop:thr-interval}(b) turns these
signs into $g_a''(0)<0$ on $(\frac12,z_-]$ and $g_a''(0)>0$ on
$[z_+,0.6712]$. The order matters. A positive interval left of $z_-$ would
give another zero, and parts (a) and (c) of Theorem~\ref{thm:threshold} would
fail. So the script prints the sign of each accepted interval with its exact
end points. Its mode \texttt{zero} fails unless the signs are $-$ before the
undecided piece and $+$ after it. Its mode \texttt{check} requires the same of
the stored records, together with the sign $-$ on every interval of the scaled
runs.
